\documentclass[Afour,sageh,times]{sagej}

\usepackage[utf8]{inputenc}
\usepackage[T1]{fontenc}
\usepackage{amsmath,amssymb}
\usepackage{graphicx}
\usepackage{booktabs}
\usepackage{multirow}
\usepackage{fontspec}
\usepackage{hyperref}
\usepackage{url}
\usepackage{xcolor}
\usepackage{listings}
\usepackage{enumitem}
\usepackage{array}
\usepackage{tabularx}
\usepackage{adjustbox}
\newfontfamily\hebrewfont{Arial}
\newcommand{\texthebrew}[1]{{\hebrewfont #1}}
\newcommand{\code}[1]{\texttt{#1}}
\setcitestyle{numbers,square}

\lstset{
  basicstyle=\ttfamily\small,
  breaklines=true,
  frame=single,
  numbers=left,
  numberstyle=\tiny\color{gray},
}

\begin{document}

\runninghead{Goldberg et al.}

\title{Distant Supervision from Cataloging Workflows: Multi-Strategy MARC-to-Wikidata Conversion for Hebrew Manuscript Collections}

\author{Alexander Goldberg\affilnum{1}, Gila Prebor\affilnum{2}, Avshalom Elmalech\affilnum{3}}

\affiliation{\affilnum{1}Department of Information Science and Applied Artificial Intelligence, Bar-Ilan University, Ramat Gan, Israel\\
\affilnum{2}Department of Information Science and Applied Artificial Intelligence, Bar-Ilan University, Ramat Gan, Israel\\
\affilnum{3}Department of Information Science and Applied Artificial Intelligence, Bar-Ilan University, Ramat Gan, Israel}

\corrauth{Alexander Goldberg, Department of Information Science and Applied Artificial Intelligence, Bar-Ilan University, Ramat Gan, Israel.}
\email{golbera3@biu.ac.il}

\begin{abstract}
Transforming library catalogs into Linked Open Data is essential for connecting isolated cultural heritage collections to the Semantic Web, yet automated conversion from MARC~21 records to Wikidata remains a largely unsolved problem---particularly for non-Latin-script collections where existing tools achieve low field coverage. We present the MHM (Mapping Hebrew Manuscripts) Pipeline, a multi-strategy system that turns Hebrew manuscript MARC records into reviewed, published Linked Open Data through three complementary contributions: \emph{authority enrichment as the primary engine of new information}, \emph{publication in three Linked-Open-Data layers}, and \emph{a curator-facing workbench that keeps researchers in the loop}.

The authority layer feeds \emph{thirteen MARC tags} (100/110/111, 600/610, 651, 700/710/711, 752, 800/810/811) into four complementary authority sources---Mazal/NLI, VIAF, Wikidata, and KIMA---and amplifies each match through \emph{VIAF cluster harvesting}, which turns a single VIAF hit into 2--4 additional international identifiers (GND, LCCN, ISNI, BnF), and \emph{Wikidata date back-fill}, which retrieves P569/P570 for SPARQL-only matches so the date guardian can vet them. The publication layer emits the same data in three forms: the full HMO research graph as RDF, a project Wikibase (\texttt{mhm-hmo.wikibase.cloud}) for browsing and editing, and a public Wikidata projection that reaches 87\% of HMO classes and averages 30.7 statements per manuscript. The workbench layer, \emph{Wikidata Studio}, surfaces every authority match, NER entity and KIMA place under a Source~$\times$~Type~$\times$~Role filter, opens a biographic-comparison dialog that aligns MARC, Mazal, VIAF and KIMA biodata side-by-side, exposes a confidence tooltip that explains why each match scored as it did, runs 14 pre-upload guards, and saves resumable review state---so every published statement is one a human approved with the evidence in view.

End-to-end evaluation on a pinned 68-record verification subset shows 97\% date coverage, 57\% genre coverage, 30.7 Wikidata statements per manuscript on average, and a $1.7\times$ increase in international identifiers per person via VIAF cluster harvesting. The three domain-specific NER models (person, provenance, contents) reach 80.4\% strict span+role F1 for role-aware person NER on a representative held-out sample, 95.9\% F1 for provenance NER, and 99.99\% F1 for contents NER; the person model also reaches 86.8\% name-only F1 and 92.6\% role accuracy when the name span is matched. The latest end-to-end evaluation finds the models' \emph{marginal} contribution beyond authority and rules to be near zero on the structured majority of records; we retain them as a safety net for the long tail of records whose entities live only in free-text notes. A distant-supervision genre classifier reaches micro-F1 0.918 ($\pm$0.003, 5-fold CV) on 8 genre classes, trained on 25{,}421 quality-filtered records with DictaBERT warm-started from the NER checkpoint and a NOTA abstention class.

The verification subset is generated deterministically by \texttt{scripts/build\_test\_subset.py}, pinned by SHA256, and used by the paper-verification harness so reported coverage and contribution numbers are reproducible. All code, models, training-data generation scripts and the Wikidata Studio application are open-source.

\end{abstract}

\begin{keywords}
Linked Open Data; Wikidata; MARC~21; named entity recognition; distant supervision; cultural heritage; Hebrew manuscripts; authority linking
\end{keywords}

\maketitle


%% ====================================================================
\section{Introduction}
\label{sec:introduction}
%% ====================================================================

\subsection{The MARC-to-LOD Challenge}

Libraries worldwide hold vast collections of cultural heritage materials cataloged in MARC~21 (Machine-Readable Cataloging), a standard developed in the 1960s that remains the dominant bibliographic record format~\cite{loc-marc}. While MARC records encode rich metadata---titles, dates, person names, physical descriptions, provenance histories---this information is trapped in a format designed for card catalog automation rather than Semantic Web interoperability.

The National Library of Israel (NLI) maintains 123,621 Hebrew manuscript MARC records, representing one of the world's largest collections of Jewish cultural heritage materials. These records contain structured fields (fixed-length codes, controlled headings) alongside unstructured note fields (MARC~500, 505, 561) where catalogers recorded detailed descriptions in Hebrew free text. The structured-unstructured duality creates both a challenge and an opportunity: rule-based conversion handles only the structured portion, while the note fields---which often contain the richest information about authorship, provenance, and contents---require natural language understanding.

Wikidata~\cite{vrandecic2014wikidata}, the collaborative knowledge base of the Wikimedia ecosystem, has emerged as the preferred Linked Open Data (LOD) target for cultural heritage institutions~\cite{van2020wikidata}. The WikiProject Manuscripts\footnote{\url{https://www.wikidata.org/wiki/Wikidata:WikiProject_Manuscripts}} defines a community-driven data model for manuscript description. However, no existing system achieves more than approximately 50\% field coverage when converting non-English MARC records to Wikidata.

\subsection{Research Gap}

Existing MARC-to-LOD conversion approaches fall into three categories, each with significant limitations:

\begin{enumerate}[leftmargin=*]
\item \textbf{Rule-based converters} (BIBFRAME~\cite{mcallister2017bibframe}, LD4L~\cite{cornell-ld4l}) handle structured MARC fields but cannot extract information from free-text notes, leaving approximately 60\% of available metadata unmapped.
\item \textbf{NER-based extractors} for cultural heritage~\cite{ehrmann2023named} exist for English and European languages but not for Hebrew bibliographic descriptions, and none integrate with Wikidata's specific property model.
\item \textbf{LLM-based approaches} (GPT-4 prompting) show promise for zero-shot entity extraction but are expensive (\$0.03--0.06/record), non-reproducible, and cannot operate offline---a critical requirement for institutions handling sensitive materials.
\end{enumerate}

No published work combines NER with rule-based extraction, multi-authority linking, and LOD enrichment in a single end-to-end pipeline that produces validated Wikidata items. Furthermore, no study has quantified the data loss inherent in MARC-to-Wikidata conversion, identifying which information is structurally untransferable.

\subsection{Contributions}

This paper makes seven contributions, organised around the three pillars announced in the abstract---authority enrichment, three-layer Linked-Open-Data publication, and a curator workbench---with the distant-supervision NER pipeline kept as a documented but secondary contribution:

\begin{enumerate}[leftmargin=*]
\item \textbf{Authority enrichment across thirteen MARC tags as the primary engine of new information.} The pipeline feeds 100/110/111, 600/610, 651, 700/710/711, 752 and 800/810/811 into four complementary authority sources (Mazal/NLI, VIAF, Wikidata, KIMA). On the pinned evaluation subset this is the layer that accounts for the majority of new statements per manuscript and converts every matched entity from an opaque string into a linked, internationally-identified item.
\item A \textbf{VIAF cluster harvesting} technique that turns a single VIAF match into 2--4 additional international authority identifiers (GND, LCCN, ISNI, BnF), yielding a $1.7\times$ increase in international identifiers per matched person and---combined with Wikidata date back-fill via P569/P570---making the date guardian effective on SPARQL-only candidates.
\item \textbf{MARC-to-LOD publication in three coordinated layers.} The same pipeline emits the full HMO research graph as RDF, loads it into a project Wikibase at \texttt{mhm-hmo.wikibase.cloud} for browsing and curator editing, and projects a public-facing view onto Wikidata that reaches 87\% of HMO classes and averages 30.7 statements per manuscript. To our knowledge this is the first time the NLI Hebrew manuscript catalog has been linked, end-to-end, to the public Semantic Web.
\item \textbf{Wikidata Studio: a curator-facing workbench that puts the researcher in control.} The workbench ships in both a cross-platform PyQt6 desktop application and a collaborative browser-based port (\texttt{mhm-pipeline-web}). Both surfaces expose every authority match, NER entity and KIMA place under a Source~$\times$~Type~$\times$~Role filter; a biographic-comparison dialog aligns MARC, Mazal, VIAF and KIMA biodata side-by-side (with Unicode NFKC normalisation, particle stripping, and side-aware highlighting); per-item override dialogs for Wikidata labels, descriptions, aliases and statements; a confidence tooltip explains which guardians fired and which cluster identifiers came along with the match; 14 pre-upload guards block malformed items; and review state is persisted so curators can pause and resume. The browser port adds shared multi-curator approval, a persistent versioning log, and encrypted per-user API keys. When the eval-agent produces a high-confidence suggested fix for a NER entity, an \textbf{Auto-fix} button appears on the relevant table row and in the detail drawer, letting the curator apply the correction in one click without pre-approving the entity. This is what makes every published statement a human-approved one and is what turns the pipeline into a research tool rather than a script.
\item A \textbf{distant-supervision methodology} that exploits the dual structure of MARC records---using structured fields as automatic labels for the same information in free-text notes---and yields three domain-specific NER models for Hebrew manuscript descriptions. The current role-aware person model reaches 80.4\% strict span+role F1 on a representative held-out benchmark sample (86.8\% name-only F1; 92.6\% role accuracy when the name span is matched), while provenance and contents NER reach 95.9\% and 99.99\% F1 respectively. The end-to-end evaluation (\S\ref{sec:evaluation}) shows their \emph{marginal} contribution beyond authority and rules is near zero on the structured majority of records; they remain in the pipeline as a safety net for the long tail in which entities only ever appear in free-text notes.
\item The first \textbf{end-to-end quantitative evaluation} of MARC-to-Wikidata conversion that measures both coverage (statements per entity) and data loss (what cannot transfer and why), together with a \textbf{five-category taxonomy of data loss} showing that the dominant remaining loss is structural---work-title strings must be reconciled to or created as Wikidata items before P1574 can be asserted---rather than technical.
\item A \textbf{complete open-source system}\footnote{\url{https://github.com/alexandergolbergwix/pipeline}} that processes 123,621 records without manual annotation, generates all training data automatically, runs the three-layer LOD publication, and ships Wikidata Studio in three deployment modes from a single codebase: a cross-platform PyQt6 desktop application, a collaborative browser-based curator workbench (FastAPI + React), and a pay-per-call Modal-hosted Stage~2 inference backend that loads all four NER and classifier models into one serverless container (\S\ref{sec:deployment}). The four model checkpoints are released as public Hugging Face Hub repositories\footnote{\url{https://huggingface.co/alexgoldberg/hebrew-manuscript-joint-ner-v2}, \url{https://huggingface.co/alexgoldberg/hebrew-manuscript-provenance-ner}, \url{https://huggingface.co/alexgoldberg/hebrew-manuscript-contents-ner}, \url{https://huggingface.co/alexgoldberg/hebrew-manuscript-genre-classifier}.} so external researchers can reuse them without redistributing weights.
\end{enumerate}


%% ====================================================================
\section{Related Work}
\label{sec:related}
%% ====================================================================

\subsection{MARC-to-Linked Data Conversion}

The Library of Congress BIBFRAME initiative~\cite{mcallister2017bibframe} defines an RDF vocabulary for bibliographic data, but BIBFRAME conversion tools (marc2bibframe2) are rule-based and limited to structured MARC fields. The LD4L (Linked Data for Libraries) project~\cite{cornell-ld4l} explored enriching library records with LOD links but focused on monograph catalogs rather than manuscripts. MARC2WIKI, a community tool for Wikidata, provides basic field mapping but achieves low coverage on non-English records due to limited date parsing and authority resolution.

The Digital Scriptorium\footnote{\url{https://digital-scriptorium.org/}} has pioneered manuscript description on Wikidata using the WikiProject Manuscripts data model, but their workflow requires substantial manual curation. Our system automates this process while maintaining compliance with the same data model.

A practitioner alternative to dedicated conversion pipelines is the widely used MarcEdit~+~OpenRefine workflow: MarcEdit provides record-level batch editing and a Linked Data Framework that embeds reconciled URIs in MARC \texttt{\$0} subfields~\cite{marcedit}, while OpenRefine reconciles tabular values against external authorities and can push statements to Wikidata through its Wikibase extension~\cite{openrefine-recon,openrefine-wikibase}. Both tools, however, operate at the record and cell level: OpenRefine's reconciliation ``happens by default through string searching''~\cite{openrefine-recon} and acts on one cell value at a time, and its Wikibase upload performs no duplicate or conflict detection against existing items~\cite{openrefine-wikibase}---precisely the failure mode our four-source cross-identifier verification and pre-upload guards are engineered against (\S\ref{sec:discussion}). Neither tool extracts entities from free-text note fields, represents typed relationships or provenance chains, or models attribution uncertainty---and the FRBRisation literature shows that entity--relationship structure can be recovered from MARC only when it is already explicitly and consistently encoded~\cite{aalberg2013}, which is rarely the case for legacy manuscript cataloging. The workflow therefore yields cleaner \emph{records}, not a queryable entity \emph{graph}: a question such as ``which manuscripts passed through the same owners' hands, and which works do they share?'' cannot be posed to a cleaned MARC file or an OpenRefine project at all, while it is a single SPARQL query over the HMO graph our pipeline produces.

\subsection{NER for Cultural Heritage}

Named Entity Recognition for historical documents has advanced significantly~\cite{ehrmann2023named}, with systems like HIPE~\cite{ehrmann2022extended} establishing benchmarks for multilingual historical NER. However, these focus on newspaper text rather than bibliographic descriptions, and none target Hebrew manuscript catalogs.

DictaBERT~\cite{shmidman2023dictabert}, a Hebrew BERT model pre-trained on diverse Hebrew corpora including historical texts, provides a strong foundation for domain-specific fine-tuning. Our prior work~\cite{goldberg2025person} demonstrated that DictaBERT can be fine-tuned for person NER on MARC note fields; this paper extends that approach to three entity types with a novel distant supervision methodology.

\subsection{Distant Supervision and Data Programming}

Distant supervision~\cite{mintz2009distant} generates training labels by aligning entity mentions in text with entries in a knowledge base. Snorkel~\cite{ratner2017snorkel} generalized this to data programming with labeling functions. Our approach differs from both: rather than using an external knowledge base, we exploit the \emph{internal dual structure} of MARC records where structured fields (100/700 for persons, 505 for contents) provide labels for the same information expressed as free text in note fields (500, 561).

\subsection{Wikidata for Cultural Heritage}

Wikidata has become a hub for cultural heritage LOD~\cite{van2020wikidata}, with initiatives like WikiProject~Books, WikiProject~Manuscripts, and Sum~of~All~Paintings linking millions of cultural objects. Feitosa et al.~\cite{feitosa2023wikidata} surveyed Wikidata usage in libraries, finding growing adoption but limited automation. Our work addresses this gap by providing a fully automated pipeline from MARC to Wikidata with property-level type validation.


%% ====================================================================
\section{System Architecture}
\label{sec:architecture}
%% ====================================================================

The MHM Pipeline consists of six stages, each producing an intermediate artifact that can be independently inspected and validated.

\subsection{Pipeline Overview}

\begin{enumerate}[leftmargin=*]
\item \textbf{MARC Parsing} (\texttt{UnifiedReader} + \texttt{field\_handlers.py}): Extracts 59 structured fields from MARC~21 records into a normalized \texttt{ExtractedData} representation, handling MARC's complex subfield structure and encoding variations.

\item \textbf{NER Extraction} (3 NER models + 2 sentence-level classifiers): Applies three DictaBERT-based NER models in parallel to different MARC fields:
  \begin{itemize}
    \item Person NER on MARC~500 (general notes): extracts person names with role classification
    \item Provenance NER on MARC~561 (ownership history): extracts owners, dates, collections
    \item Contents NER on MARC~505 (contents): extracts work titles, folio references, work authors
  \end{itemize}
  A complementary classifier (genre, Strategy~F) runs alongside the NER models in the same Stage~2 worker, sharing the DictaBERT encoder.

\item \textbf{Authority Resolution}: Matches names against four authority sources: Mazal/NLI (local SQLite, $\sim$2.5~million records), VIAF (international, SRU API), KIMA (Hebrew place gazetteer), and Wikidata (SPARQL, public). The set of MARC name- and place-bearing fields consulted at this stage was expanded in v1.10 (\texttt{CLAUDE.md}~rule~49) from the six original tags (100/110/111 main entry, 700/710/711 added entry) to thirteen: \emph{subject} headings 600/610/651, \emph{hierarchical place name} 752, and \emph{series added entries} 800/810/811 now feed the matcher alongside the original six. Subject persons (MARC~600) take a dedicated \texttt{role="subject"} that suppresses the death-side date guard --- a manuscript can be \emph{about} someone who died centuries before its production. VIAF cluster harvesting extracts cross-referenced GND, LCCN, ISNI, and BnF identifiers.

\item \textbf{RDF Construction}: Maps curator-reviewed authority-enriched records to the Hebrew Manuscripts Ontology (HMO), producing the canonical \texttt{output.ttl} graph compliant with FRBRoo and CIDOC-CRM. This graph preserves the full scholarly model, including codicological units, expressions, transmission witnesses, text traditions, and evidential structures that are not always suitable for Wikidata.

\item \textbf{SHACL Validation}: Validates the RDF graph against SHACL shapes encoding cardinality constraints, value type requirements, and mandatory properties.

\item \textbf{Wikidata Projection and Upload}: Projects the validated HMO graph into WikibaseIntegrator item representations and uploads or exports them via the Wikidata API / QuickStatements with OAuth~2.0 authentication, claim-level diffing (no duplicate statements), and batch rate limiting. The same stage also emits a projection-coverage report and an offline project-Wikibase draft export. The latter is designed as a full-HMO deployment draft: each typed HMO node becomes a project-Wikibase draft entity and RDF statements are preserved as local statement drafts, while Wikidata receives only the conservative public projection.
\end{enumerate}

Every stage exposes a substep telemetry channel (\texttt{substep~=~pyqtSignal(str)} on the \texttt{StageWorker} base class) so the GUI's unified \texttt{DynamicProgressBar} surfaces the running operation (e.g.\ ``Stage 3.1~--~Mazal lookup ($i/n$)'') alongside percentage and rolling-window ETA. All nine workers --- the six pipeline stages plus the three preview/validate variants --- emit through the same channel, so the user sees one consistent progress convention regardless of which stage is running.

\subsection{Deployment Architecture: Desktop, Browser, and Pay-per-Call Inference}
\label{sec:deployment}

The same pipeline runs in three deployment modes, each appropriate for a different operational context. The model checkpoints, post-filters, authority matchers, and Wikidata mapping are byte-identical across modes; only the inference \emph{host} differs.

\paragraph{Desktop application.} A cross-platform PyQt6 executable (macOS DMG, Windows installer) bundles the four \texttt{.pt} checkpoints (~2.5\,GB total) and a DictaBERT base alongside the application. Stage~2 inference runs in-process via PyTorch with MPS$\rightarrow$CUDA$\rightarrow$CPU device fallthrough. This is the canonical research deployment: full offline operation, no per-call cost, and the lowest latency (~50\,ms per record on Apple Silicon). The bundled-models lockstep contract is enforced by a cross-platform parity test that fails the build if either installer drops a checkpoint.

\paragraph{Collaborative web port.} The same pipeline backend is exposed as a FastAPI service with a React/TypeScript frontend at \texttt{mhm-pipeline-web} (project repository), deployable to standard PaaS environments. The browser surface mirrors the desktop's curator workflow: per-row authority approval, agentic AI verification, and a shared bidirectional approval store so multiple curators can review the same run live. Every curator surface supports per-field manual override---NER entities (\texttt{override\_text/type/role}), authority matches, Wikidata Studio item labels/descriptions/statements, and raw MARC JSON---with a per-entity versioning log (RFC~6902 patches) and project export bundles so research teams can audit and replay decisions. Encrypted per-user secrets (Gemini key, Wikidata token, Wikibase Cloud bot password) are stored under a password-derived KEK (zero-knowledge recovery).

\paragraph{Pay-per-call inference backend.} To support web deployments without bundling model weights into the application slug, all four Stage~2 models (joint Person NER, Provenance NER, Contents NER, Genre Classifier) are also published as public Hugging Face Hub repositories\footnote{\url{https://huggingface.co/alexgoldberg/hebrew-manuscript-joint-ner-v2}, \url{https://huggingface.co/alexgoldberg/hebrew-manuscript-provenance-ner}, \url{https://huggingface.co/alexgoldberg/hebrew-manuscript-contents-ner}, \url{https://huggingface.co/alexgoldberg/hebrew-manuscript-genre-classifier}.} and bundled together in a single serverless container hosted on Modal.\footnote{\url{https://modal.com/}.} A single \texttt{MhmNer} class loads all four pipelines into one container at warm-up (sharing the DictaBERT encoder across the three derived NER models), exposes a single \texttt{POST /extract} endpoint that returns all four model outputs in one round-trip per record, and scales to zero after 5~minutes of inactivity. Cold start: ~30--60\,s for the four-model load; warm latency: ~300\,ms per record on a 2-CPU container. Billed per-second of container time, this works out to approximately \$0.00006 per record (well under \$1 for a 10{,}000-record corpus). Modal's free-tier credit covers typical research workloads entirely. The web FastAPI backend selects between \emph{local} (in-process torch), \emph{HuggingFace Inference Providers}, and \emph{Modal} extraction backends through an \texttt{EXTRACTION\_MODE} environment variable; the choice is transparent to the rest of the pipeline. A shared cross-user inference cache routes every external NER, classifier, and authority call through a two-tier lookup: an in-process Redis L1 tier (sub-millisecond, content-addressed by SHA-256 of canonical query JSON) backed by a Postgres L2 tier (durable, 30-day TTL for authority results; no expiry for NER/classifier). The first team member to resolve a candidate populates both tiers; every subsequent request warm-hits Redis without touching Postgres.

The three deployment modes preserve the same Stage~3--6 logic and the same safety guards. Communication with external inference services is HTTP-only; no Python imports cross the FastAPI backend $\leftrightarrow$ Modal trust boundary, mirroring the equivalent boundary between the desktop pipeline and the eval-agent subprocess (Rule~48 in the project operating manual).

\subsection{Wikidata Property Mapping}

The pipeline maps to 37 Wikidata properties following the WikiProject Manuscripts data model. Table~\ref{tab:manuscript-props} shows manuscript properties; Table~\ref{tab:person-props} shows person properties. In the current architecture, these properties constitute the \emph{public projection} of the richer HMO graph: a machine-readable \texttt{wikidata\_projection\_coverage.json} report records, for each HMO class present in \texttt{output.ttl}, whether it is projected as a Wikidata item, summarized through claims or qualifiers, or intentionally retained only in HMO / project Wikibase.

\begin{table}[t]
\caption{Wikidata properties for manuscript entities (pinned 68-record subset, 2026-07 build). P50 (author) is intentionally absent: the community constraint on Property:P50 requires the exemplar chain manuscript $\to$ P1574 $\to$ work $\to$ P50, so authorship is asserted on work items, never directly on manuscripts.}
\label{tab:manuscript-props}
\centering
\footnotesize
\setlength{\tabcolsep}{3.5pt}
\begin{tabular}{llll}
\toprule
\textbf{Property} & \textbf{Label} & \textbf{Source} & \textbf{Coverage} \\
\midrule
P31  & instance of      & Rule       & 100\% \\
P195 & collection       & Hardcoded  & 100\% \\
P17  & country          & Hardcoded  & 100\% \\
P131 & located in       & Hardcoded  & 100\% \\
P217 & inventory no.    & MARC 852   & 96\% \\
P3959& NLI system no.   & MARC 001   & 100\% \\
P1476& title            & MARC 245   & 94\%  \\
P407 & language         & MARC 008   & 100\% \\
P282 & writing system   & Rule       & 100\% \\
P571 & inception        & MARC 008   & 97\%  \\
P136 & genre            & MARC 655 + ML & 57\%  \\
P973 & described at URL & NLI catalog & 100\% \\
P2635& number of parts  & CU count   & 3\%  \\
P921 & main subject     & MARC 630   & 76\%  \\
P7153& significant place& KIMA       & 7\%  \\
P1071& location created & KIMA       & 7\%  \\
P1684& inscription      & Colophon   & 6\%  \\
P11603&transcribed by   & NER+Auth   & 1\%  \\
P127 & owned by         & Prov.\ NER & 21\%  \\
P6216& copyright        & Rule       & 91\% \\
P5008& on focus list    & Hardcoded  & 100\% \\
\bottomrule
\end{tabular}
\end{table}

\begin{table}[t]
\caption{Wikidata properties for person entities (verified on 66 persons in the pinned subset build; the 2026-06 authority-hardening pass---homonym abstention and a fail-closed no-identifier gate---builds fewer but better-grounded person items).}
\label{tab:person-props}
\centering
\small
\begin{tabular}{lllr}
\toprule
\textbf{Property} & \textbf{Label} & \textbf{Strategy} & \textbf{Count} \\
\midrule
P31   & instance of    & Rule            & 66 \\
P1412 & language       & MARC 008/041    & 78 \\
P1559 & native name    & Hebrew name     & 64 \\
P106  & occupation     & Role mapping    & 63 \\
P214  & VIAF ID        & VIAF search     & 61 \\
P213  & ISNI           & Cluster harvest & 39 \\
P244  & LCCN           & Cluster harvest & 37 \\
P227  & GND ID         & Cluster harvest & 26  \\
\midrule
\multicolumn{3}{l}{\textbf{Average statements per person}} & \textbf{6.6} \\
\bottomrule
\end{tabular}
\end{table}


%% ====================================================================
\section{Multi-Strategy Extraction Methodology}
\label{sec:methodology}
%% ====================================================================

\subsection{Strategy A: Distant Supervision NER}
\label{sec:distant-supervision}

The core methodological contribution is the exploitation of MARC's dual structure for distant supervision. MARC records contain the same information twice: once in structured fields with controlled vocabulary (e.g., field~100\$a for author headings) and once in free-text notes written by catalogers (e.g., field~500\$a). This creates a natural alignment between structured labels and unstructured text without any manual annotation.

\begin{table*}[t]
\caption{Training regime for the five neural models. All five use DictaBERT transformer encoders rather than feature-engineered classical models such as random forests.}
\label{tab:model-training-regime}
\centering
\small
\begin{tabularx}{\textwidth}{p{0.21\textwidth}p{0.27\textwidth}p{0.24\textwidth}p{0.18\textwidth}}
\toprule
\textbf{Model} & \textbf{Encoder update} & \textbf{Task head} & \textbf{Training type} \\
\midrule
Person NER + role & DictaBERT fully fine-tuned; all transformer layers are updated. & Role-aware token-level BIO head whose labels encode both span and role (AUTHOR, TRANSCRIBER, OWNER, CENSOR, TRANSLATOR, COMMENTATOR). & Full fine-tuning \\
Provenance NER & DictaBERT fully fine-tuned; all transformer layers are updated. & Token-level BIO head for OWNER, DATE, and COLLECTION spans, trained from scratch. & Full fine-tuning \\
Contents NER & DictaBERT fully fine-tuned; all transformer layers are updated. & Token-level BIO head for WORK, FOLIO, and WORK\_AUTHOR spans, trained from scratch. & Full fine-tuning \\
Genre classifier & Lower DictaBERT layers are frozen while the upper layers adapt with a small learning rate. & Multi-label sigmoid head over 8 retained genres plus NOTA, trained from scratch. & Partial fine-tuning \\
\bottomrule
\end{tabularx}
\end{table*}

Fine-tuning here means that the 110M-parameter DictaBERT encoder, pre-trained for Hebrew language modelling, is used as the starting point and then adapted to manuscript-catalog tasks. For the three token-level NER models, all encoder layers are updated together with a new task head. For the two classifiers, we use parameter-efficient partial fine-tuning: lower layers that encode general Hebrew morphology and syntax remain fixed, while the upper layers and the newly initialised sigmoid classification head are trained. This follows the same conservative adaptation principle as adapter or LoRA-style methods---preserve most pre-trained linguistic knowledge and update only the task-specific capacity---but no LoRA adapters are used; the resulting checkpoints are ordinary DictaBERT classifier checkpoints that can be loaded directly by PyTorch in the desktop application, by the FastAPI backend when configured in local mode, or by the Modal serverless container that vendors the same loader code (\S\ref{sec:deployment}).

\subsubsection{Model 1: Role-aware Person NER (80.4\% strict span+role F1)}

\textbf{Training data generation.} For each MARC record, person names and roles from structured fields (100\$a, 600\$a, 700\$a, 800\$a) are matched against occurrences in note fields (500\$a, 520\$a, 545\$a, 546\$a, 561\$a, 957\$a) using exact and fuzzy string matching. When a name appears in the note text, the corresponding tokens receive role-aware BIO labels such as \texttt{B-TRANSCRIBER}, \texttt{I-OWNER}, or \texttt{B-AUTHOR}; all other tokens receive the O label. The current v3 dataset contains 6,823 training, 911 validation, and 904 held-out test samples from the NLI catalog.

\textbf{Architecture.} DictaBERT~\cite{shmidman2023dictabert} (110M parameters) with a single role-aware token-classification head. Trained as a full fine-tune and deployed as \texttt{alexgoldberg/hebrew-manuscript-joint-ner-v2} / \texttt{models/hebrew-manuscript-joint-ner-v2}.

\textbf{Role classification.} Role assignment is now learned directly by the BIO label sequence rather than delegated to a separate keyword heuristic. On the representative 100-record held-out benchmark sample, exact matching of both name span and role yields 80.4\% F1 (P~=~79.3, R~=~81.5); name extraction alone reaches 86.8\% F1, and the role is correct for 92.6\% of matched names.

\textbf{Confidence signals.} Each emitted person span carries two distinct confidence values: a legacy \texttt{confidence} field derived from the keyword classifier (bimodal at 0.60 and 0.85), preserved for backward compatibility with downstream Stage~3 thresholds, and a \texttt{model\_confidence} field reporting the mean softmax probability across the entity's tokens, which the GUI auto-approve gate consumes as an orthogonal signal.

\textbf{Why this model is needed and what it adds.} Structured MARC headings capture only the explicit main and added entries; catalogers mention many additional people --- authors of embedded works, translators, commentators, and scribes --- only inside MARC~500 note prose. One training example comes from record~001~990001253330205171, whose note segment \texthebrew{... תפסיר י' דברות מאת אלעזר בן אליעזר} yields an $X$ token sequence containing \texthebrew{אלעזר בן אליעזר} and a $Y$ sequence that labels those three tokens as \texttt{B-PERSON / I-PERSON / I-PERSON}. On the pinned 68-record verification subset, every record contains note text, so the potential reach of Person NER is 68/68 records. Re-running Stage~2 with the current code on that subset emits 59 person-role spans across 36 records (45~AUTHOR, 8~TRANSCRIBER, 4~TRANSLATOR, 1~COMMENTATOR, 1~OWNER), feeding the downstream P50 and P11603 claim paths.

\begin{figure}[t]
\centering
\footnotesize
\fbox{\parbox{0.92\linewidth}{\textbf{MARC source}\\Record 001: 990001253330205171\\MARC 500 snippet: \texthebrew{... תפסיר י' דברות מאת אלעזר בן אליעזר}}}

\vspace{0.35em}

$\Downarrow$

\vspace{0.35em}

\fbox{\parbox{0.92\linewidth}{\textbf{$X$ fed to the model}\\Token sequence from the note segment: \texthebrew{... | מאת | אלעזר | בן | אליעזר}}}

\vspace{0.35em}

$\Downarrow$

\vspace{0.35em}

\fbox{\parbox{0.92\linewidth}{\textbf{$Y$ created for training}\\BIO sequence: \texttt{... | O | B-PERSON | I-PERSON | I-PERSON}}}
\caption{Model~1 (Person NER): distant-supervision example. Structured person anchors are projected back into MARC~500 note text, producing a token sequence $X$ and a BIO label sequence $Y$.}
\label{fig:person-ner-real-example}
\end{figure}

\subsubsection{Model 2: Provenance NER (95.9\% F1)}

\textbf{Training data generation.} Three distant supervision strategies target MARC~561 (ownership/custodial history):
\begin{enumerate}[leftmargin=*]
\item \emph{Structured cross-reference}: Owner names from field~561 are aligned with authority headings in 100/700 fields of the same record.
\item \emph{Hebrew pattern matching}: Regular expressions identify ownership formulae (\texthebrew{ציון בעלים}, \texthebrew{אוסף}) and extract adjacent entity spans.
\item \emph{Latin censor identification}: Names following Latin censor attestation patterns (``Visto per mi...'') are labeled as PERSON entities.
\end{enumerate}

Training on 12,100 samples (28.4\% augmented with multi-entity examples) yielded 95.91\% F1 (best fold: 96.17\%) for three entity types: OWNER, DATE, COLLECTION.

\textbf{Encoder comparison.} As a post-hoc architecture check, we retrained the provenance task on the same 12,100-sample dataset with \texttt{dicta-il/neodictabert-bilingual} under the same 5-fold split. This experimental run reached a mean F1 of 94.65\% ($\pm$0.34; best fold: 95.14\%), compared with 95.91\% ($\pm$0.15; best fold: 96.17\%) for the production DictaBERT checkpoint. Because the NeoDictaBERT run used a smaller micro-batch and a shorter training budget for MPS stability, we treat it as an exploratory comparison rather than a definitive rejection of the newer encoder. Nevertheless, the result supports the current engineering choice: for this domain-specific sequence-labelling task, the smaller DictaBERT checkpoint remains both more accurate and more deployable in the desktop pipeline.

\textbf{Why this model is needed and what it adds.} Ownership history is almost never stored as a clean one-field/one-owner structure. Instead, MARC~561 mixes owners, dates, collections, acquisition formulas, and long documentary inscriptions in a single prose field. In record~001~990001253170205171, the segment \texthebrew{ציון בעלים: "מקנת כספי הצעיר יעקב בכ"ר שלמה הכהן"} becomes the $X$ token sequence, while the $Y$ sequence marks the owner phrase beginning at \texthebrew{מקנת} as a contiguous OWNER span. On the pinned subset, 36/68 records contain direct MARC~561 provenance text, which defines the model's core reachable gap. The current Stage~2 rerun emits 13 provenance entities across 15 records (12~OWNER, 4~DATE, 1~COLLECTION); in the cached end-to-end subset run, this ultimately yields P127 claims on 14/68 manuscripts.

\begin{figure}[t]
\centering
\footnotesize
\fbox{\parbox{0.92\linewidth}{\textbf{MARC source}\\Record 001: 990001253170205171\\MARC 561 segment: \texthebrew{ציון בעלים: "מקנת כספי הצעיר יעקב בכ"ר שלמה הכהן"}}}

\vspace{0.35em}

$\Downarrow$

\vspace{0.35em}

\fbox{\parbox{0.92\linewidth}{\textbf{$X$ fed to the model}\\Segment tokenization: \texthebrew{ציון | בעלים | מקנת | כספי | הצעיר | יעקב | בכ"ר | שלמה | הכהן}}}

\vspace{0.35em}

$\Downarrow$

\vspace{0.35em}

\fbox{\parbox{0.92\linewidth}{\textbf{$Y$ created for training}\\BIO sequence: \texttt{O | O | B-OWNER | I-OWNER | I-OWNER | I-OWNER | I-OWNER | I-OWNER | I-OWNER}}}
\caption{Model~2 (Provenance NER): MARC~561 training example. Ownership-note formulas and aligned headings produce OWNER/DATE/COLLECTION silver labels.}
\label{fig:provenance-ner-real-example}
\end{figure}

\subsubsection{Model 3: Contents NER (99.99\% F1)}

MARC~505 (contents notes) follows a highly regular semi-structured format: numbered entries with work titles, folio references, and sometimes author names. Rule-based parsing of 505 fields generates training data with near-perfect labels, explaining the exceptional F1 score. Three entity types: WORK, FOLIO, WORK\_AUTHOR. Trained on 10,000 samples.

\textbf{Why this model is needed and what it adds.} The contents note packs several distinct targets into a single string: folio range, embedded work title, and sometimes the work's author. In record~001~990001253170205171, the string \texthebrew{דף 1א-176ב: ויטל, חיים בן יוסף: אוצרות חיים} becomes the $X$ sequence, and the $Y$ labels recover one FOLIO span, one WORK\_AUTHOR span, and one WORK span. On the pinned subset only 2/68 records contain MARC~505 contents notes, but those records are dense. The current Stage~2 rerun emits 93 contents entities on those two records (56~FOLIO, 35~WORK, 2~WORK\_AUTHOR), supplying the structured evidence that the downstream work-item builder uses to produce 126~P1574 exemplar-of claims in the cached end-to-end subset run.

\begin{figure}[t]
\centering
\footnotesize
\fbox{\parbox{0.92\linewidth}{\textbf{MARC source}\\Record 001: 990001253170205171\\MARC 505 entry: \texthebrew{דף 1א-176ב: ויטל, חיים בן יוסף: אוצרות חיים}}}

\vspace{0.35em}

$\Downarrow$

\vspace{0.35em}

\fbox{\parbox{0.92\linewidth}{\textbf{$X$ fed to the model}\\Token sequence: \texthebrew{דף | 1א-176ב | ויטל | חיים | בן | יוסף | אוצרות | חיים}}}

\vspace{0.35em}

$\Downarrow$

\vspace{0.35em}

\fbox{\parbox{0.92\linewidth}{\textbf{$Y$ created for training}\\BIO sequence: \texttt{B-FOLIO | I-FOLIO | B-WORK\_AUTHOR | ... | B-WORK | I-WORK}}}
\caption{Model~3 (Contents NER): MARC~505 training example. Semi-structured contents entries are converted into token-level labels for FOLIO, WORK\_AUTHOR, and WORK.}
\label{fig:contents-ner-real-example}
\end{figure}

\subsubsection{Post-extraction filtering}
\label{sec:ner-post-filters}

The three NER models emit raw spans that are then routed through four deterministic post-filters in \texttt{converter/authority/ner\_post\_filters.py} before reaching Stage~3. (i)~A WORK\_AUTHOR span whose payload matches a folio-shape pattern (digits followed by a Hebrew side letter, e.g.\ \texthebrew{133ב~:}) is re-typed to FOLIO and stamped with a \texttt{retyped\_from} marker. (ii)~A COLLECTION span whose surrounding context lacks an institutional marker (\texthebrew{אוסף} ``collection'', ``Library'', ``ms'') and whose surname matches a curated catalog-citation list is routed to a separate \texttt{record["catalog\_references"]} channel rather than treated as an institutional collection; this prevents catalog citations like \texthebrew{מ' גסטר} (``M.\ Gaster'') from polluting P195 (collection) statements. (iii)~OWNER spans longer than 80 characters are routed to \texttt{record["provenance\_inscriptions"]} where they feed P7535 (description) narratively, rather than being collapsed into a single P127 (owned~by) statement; bills of sale and multi-line ownership inscriptions are documentary objects, not single-owner attestations. (iv)~Person spans matching curated topic-keyword denylists (e.g.\ \texthebrew{קבלה} ``kabbalah'', \texthebrew{משיח} ``messiah'', \texthebrew{קולופון} ``colophon''), all-caps ASCII fragments, or sequences with insufficient Hebrew letter count are dropped as model hallucinations. On the 68-record audit corpus, the four filters together eliminated 22 WORK\_AUTHOR/folio confusions, 5 catalog-citation false positives, 5 bill-of-sale paragraphs misrouted to P127, and 16 person-hallucination spans, with no impact on true-positive person, owner, or work entities.

\subsection{Strategy B: Rule-Based Pattern Extraction}

\subsubsection{Hebrew Date Parsing}

Hebrew manuscript dates present unique challenges. NLI catalogs encode dates in three formats: structured years (``1495''), English century descriptions (``16th century''), and Hebrew century strings (\texthebrew{מאה ט"ז} = 16th century = ca.~1550). Our parser handles all three formats, including Hebrew ordinal-to-integer conversion via a mapping of 23 entries (21 unique centuries, with two duplicate spellings for the special-case 15 and 16). Coverage improved from 22\% to 97\% after implementing Hebrew century parsing.

\subsubsection{Genre and Subject QID Mapping}

We created comprehensive mapping dictionaries: 50~genre terms to Wikidata QIDs (covering 100\% of NLI's genre vocabulary), 37~LCSH subject headings, 13~Bible books, and 14~Talmud Bavli tractates. All mappings are centralized in a single module (\texttt{property\_mapping.py}) for maintainability.

\subsection{Strategy C: Multi-Authority Linking}

Four authority sources provide entity resolution:

\textbf{Mazal/NLI} (local SQLite, $\sim$2.5~million records): Fuzzy matching of person, place, work, and corporate body names against the NLI authority file. Provides NLI~J9U identifiers (P8189) and biographical dates.

\textbf{VIAF} (public API, SRU/CQL): Name search with JSON content negotiation. Rate-limited to 2~requests/second. Provides VIAF~URIs (P214).

\textbf{KIMA} (Hebrew place gazetteer, SQLite): Maps Hebrew historical place names to Wikidata QIDs via variant name matching and stores the full KIMA index row per match---Hebrew and romanized preferred names, latitude/longitude, and GeoNames identifiers---so downstream coordinate claims (P7153, P625) and the Match Detail dialog's location panel can surface geographic evidence without a second lookup. Provides location of creation claims (P1071).

\textbf{Wikidata} (SPARQL, public): A fourth authority source resolves entities directly against Wikidata using three search modes --- (i)~identifier triangulation, when a candidate carries P214 / P8189 / P244 / P227 / P213; (ii)~Hebrew-label search ordered by ascending QID number, preferring canonical low-numbered entries (e.g.\ Q189564 for Rashi) over pipeline-created Q139xxx duplicates; and (iii)~Latin-name fallback. The matcher's \texttt{find\_viaf\_by\_qid} method reads \texttt{wdt:P214} off any known QID; this is invoked as a backfill step when a Mazal authority match resolved to a Wikidata QID without a VIAF URI, increasing VIAF coverage from 13\% to 33\% on the 68-record audit corpus and unlocking the cluster-harvesting amplification described in Strategy~D. A companion method \texttt{find\_dates\_by\_qid} (added in v1.10, \texttt{CLAUDE.md}~rule~49) queries \texttt{wdt:P569} (date of birth) and \texttt{wdt:P570} (date of death) in a single \texttt{OPTIONAL}-block SPARQL query so that the role-aware date-conflict guard fires on Wikidata-only candidates as well, not only on those resolved by VIAF or Mazal. To survive community-edited disagreement, the method abstains on any property whose SPARQL response carries two or more distinct year values, treating that field as \texttt{None}. The resulting birth/death years feed a three-tier precedence chain \texttt{viaf\_birth ${\,\vee\,}$ mazal\_birth ${\,\vee\,}$ wd\_birth} (and analogously for death) into the existing guard; on the 68-record audit corpus this added biographical-date coverage on roughly one-third of the persons resolved purely by Wikidata SPARQL, with no impact on candidates already dated by VIAF or Mazal.

\textbf{Date-conflict guard} (\texttt{stage3\_guards.evaluate\_date\_conflict}). Independent of which authority supplied the dates, every person match is checked against the manuscript's catalogued production year. Roles in \texttt{PHYSICAL\_PRODUCTION\_ROLES} (scribe, transcriber, copyist) get the death check (rejected if the candidate died more than 80~years before the manuscript was made); \texttt{TEXTUAL\_AUTHORSHIP\_ROLES} (author, translator, commentator, editor) get only the birth check because Hebrew manuscripts routinely copy medieval authors centuries after their death (Maimonides~d.~1204 in 17th-century copies, Rashi~d.~1105, etc.); and the new \texttt{role="subject"} introduced for MARC~600 subject persons gets only the birth check because a manuscript can be \emph{about} a long-deceased figure but cannot be about someone born after it was made. A 5-year buffer (\texttt{DATE\_BIRTH\_BUFFER\_YEARS}) absorbs typical round-year cataloguing imprecision.

Each match record carries a \texttt{source} field --- one of \texttt{mazal}, \texttt{viaf}, \texttt{wikidata}, \texttt{cross\_source}, or \texttt{marc\_only} --- and a \texttt{source\_count} integer in $[0, 3]$ recording how many independent identifier sources agreed on the resolution. Matches with \texttt{source\_count} $\geq 2$ also carry a \texttt{sources} array enumerating the agreeing sources (e.g.\ \texttt{["mazal", "wikidata"]}). The \texttt{score\_confidence} function consumes these signals plus VIAF over-merge detection and cross-source conflict probes to assign one of \texttt{high}, \texttt{medium}, or \texttt{low}; only \texttt{high} survives the GUI's 0.95 auto-approve threshold, while \texttt{medium} and \texttt{low} require explicit curator review.

\subsection{Strategy D: VIAF Cluster Harvesting}

When a person matches a VIAF cluster, we fetch the full cluster JSON from \texttt{https://viaf.org/viaf/\{id\}} (with \texttt{Accept: application/json} header) and extract cross-referenced authority identifiers. The VIAF cluster aggregates records from 40+ national libraries; a single match thus yields up to four additional external identifiers. In the browser port, \texttt{payload.cluster\_ids} on each \texttt{AuthorityMatch} row stores the harvested GND/LCCN/ISNI/BnF/J9U map from the VIAF metadata response (an early web-port regression had left this field empty despite successful VIAF hits; the fix ensures the Sources tab and Wikidata item builder see the same cluster the desktop pipeline emits):

\begin{itemize}[leftmargin=*]
\item \textbf{P213} (ISNI): International Standard Name Identifier --- 39 persons (59\%)
\item \textbf{P244} (LCCN): Library of Congress --- 37 persons (56\%)
\item \textbf{P227} (GND): Deutsche Nationalbibliothek --- 26 persons (39\%)
\item \textbf{P268} (BnF): Biblioth\`{e}que nationale de France --- none in the current subset build
\end{itemize}

This technique transforms a single VIAF URI into a web of cross-referenced international authority links, significantly increasing each person entity's LOD connectivity at minimal additional API cost (one HTTP request per person).

\subsection{Strategy E: Property-Level Validation}

Wikidata properties have strict value type constraints. Our validator checks every generated claim against the expected data type:
\begin{itemize}[leftmargin=*]
\item P8189 and P214 require \texttt{external-id}, not \texttt{string}
\item P5816 (condition) requires item QIDs, not free text --- claims with string values are skipped
\item P527 (has parts) requires item QIDs --- work titles without Wikidata items are omitted
\item P195 (collection) requires item QIDs --- NER-extracted collection names without reconciled QIDs are skipped
\end{itemize}

This validation prevents upload errors and ensures WikiProject Manuscripts Data Model compliance.

\subsection{Strategy F: Distant-Supervision Genre Classifier}

31\% of NLI Hebrew manuscript records have no MARC~655 genre/form headings, leaving P136 (genre) empty after rule-based QID mapping. To fill this gap we train a multi-label sequence classifier via distant supervision: records that \emph{do} carry MARC~655 labels provide (title + context window, genre) training pairs, and the trained model is applied to the 31\% without labels.

\textbf{Data and filtering.} The full Ktiv catalog export (123,621 records) is used as the training corpus, augmented with 897 intensively annotated records. A record is included only if three conditions hold: (1)~MARC~245 (title), MARC~500 (general notes), and MARC~655 (genre) are all present; (2)~at least one genre keyword appears in MARC~500 notes as a \emph{whole token}; and (3)~the genre maps to a key in \texttt{GENRE\_TO\_QID}. Condition~(2) is a quality gate: it confirms the cataloger explicitly described the genre in free text, ensuring a genuine textual signal exists for learning. Whole-token matching (space-boundary for Hebrew, \texttt{\textbackslash b} for ASCII) prevents spurious substring matches (e.g., \texthebrew{שיר} matching \texthebrew{שירות}). A pre-extraction step saves the 26,758 qualifying records to a compact TSV (\texttt{genre\_samples.tsv}), making all subsequent training runs fast without re-scanning the full catalog.

Two additional quality filters are applied: (i)~\emph{Literature (Miscellaneous, in manuscript)} is excluded because its keywords (\texthebrew{קובץ}, \texthebrew{אסופה}, ``collection'') are generic enough to appear across all genres, making its 2,170 samples predominantly false positives; (ii)~genre classes with fewer than 100 keyword-matched training examples are discarded (Business records: 75, Family records: 63, Biographies: 34, Negotiable instruments: 20), as they contribute more label noise than learning signal. After both filters, the training set contains 25,421 records across 8 high-confidence classes plus 1,629 NOTA examples.

\textbf{Keyword synonym dictionary.} Since MARC~655 headings are in English while MARC~500 notes are in Hebrew, direct string overlap yields near-zero matches. We therefore compiled \texttt{GENRE\_KEYWORDS}: a dictionary mapping each of the 20 genre classes to all known Hebrew morphological forms and English variants. For example, \emph{Legislation (Jewish law)} maps to \{תקנה, תקנות, תקנן, תוקן, חוק, חוקים, פסק, פסקי דין, גזרה, גזרות, חרם, חרמות, צו, צווים, takkanot, decree, ordinance, statute\}.

\textbf{Input (X).} Rather than using the entire notes field, the model input is the MARC~245 title concatenated with a \emph{3-sentence context window} from MARC~500: the sentence immediately before the keyword match, the matching sentence itself, and the sentence immediately after. This localises the input to the exact textual evidence the cataloger wrote about that genre. Sequences are truncated to 64 tokens (matched to the inference window).

\textbf{Architecture.} The encoder is \texttt{dicta-il/dictabert}~\cite{shmidman2023dictabert} warm-started from the provenance NER checkpoint (already domain-adapted on 12,100 Hebrew manuscript samples via our earlier distant supervision pipeline). The lower transformer layers are frozen --- they encode generic Hebrew morphology and syntax already learned during pre-training and NER fine-tuning --- while the upper layers and the classification head are fine-tuned. This partial fine-tuning allows the encoder to adapt its higher-level representations to the genre domain without catastrophic forgetting of Hebrew pretraining. The [CLS] representation is passed through Dropout(0.3) and a linear layer to logits over 9 classes: the 8 retained genres plus an explicit \emph{None-of-the-Above} (NOTA) class, allowing the model to abstain when no genre applies. Because this is a multi-label task, logits are passed through independent sigmoid functions rather than a softmax: a manuscript may legitimately be both \emph{Poetry} and \emph{Illustrated works}. The sigmoid head is trained from scratch; it is a newly initialised linear projection from DictaBERT's 768-dimensional [CLS] representation to the 9 genre logits. Focal loss~\cite{lin2017focal} with $\gamma=2$ and per-class positive weights $w_c = n_{\text{neg},c}/n_{\text{pos},c}$ is used instead of standard BCE, downweighting easy majority-class examples and focusing training on the harder, rarer classes.

\textbf{Training.} AdamW with differential learning rates: top 2 encoder layers at $2 \times 10^{-6}$ (10$\times$ lower than the head at $2 \times 10^{-5}$) to prevent catastrophic forgetting. Batch size 64, max sequence length 64 tokens. 5-fold stratified cross-validation, up to 30 epochs per fold, early stopping (patience~=~5), linear warmup over 10\% of steps. The best-fold checkpoint is saved. The classification threshold is tuned on each fold's validation set by scanning $[0.2, 0.8]$ in steps of 0.05 to maximise micro-F1; the optimal threshold is stored in the checkpoint and used at inference. Total training time on Apple M4 Pro (MPS): approximately 1.5 hours.

\textbf{Output and integration.} At inference, the classifier is applied to any record lacking MARC~655 genres. Because the model was trained on short 3-sentence context windows (typically under 64 tokens) but the inference input combines the title with up to three full MARC~500 notes (which can be substantially longer), a \emph{sliding window} strategy is used: the full token sequence is divided into overlapping 64-token windows with 50\% stride, each window is scored independently, and the sigmoid probabilities are averaged across all windows before thresholding. This prevents information loss from naive truncation. Genres predicted above the tuned threshold are added as P136 statements with epistemic qualifiers: P1480 = Q18122778 (\emph{presumably}) and a P887 reference citing Q2539 (\emph{machine learning}), distinguishing classifier-inferred claims from authority-sourced data. When the model predicts NOTA — or no genre exceeds the threshold — the classifier returns \emph{``other''} and no P136 claim is written. The model is loaded lazily and degrades gracefully when absent.

\textbf{Why this model is needed and what it adds.} The genre classifier exists specifically for the records whose genre is \emph{missing}, not for those already labelled in MARC~655. One training example is record~001~990001273190205171: MARC~245 contributes the title \texthebrew{סדור מנהג אפגניסטן}, MARC~500 contributes the local context \texthebrew{כולל פיוט מאת פנחס חכירי}, and MARC~655 provides the distant-supervision target \texttt{Piyyutim}. Thus $X$ is title~$+$~local note context and $Y$ is a multi-label genre vector whose positive bit is \texttt{Piyyutim}. On the pinned 68-record subset, 40 records have no raw genres at all; this is the exact missing-data gap addressed by the model. The current Stage~2 rerun emits 30 genre predictions across 28 records, and 11 of those records had no raw MARC~655 genre, expanding the subset from 28 records with native genre data to 39 records with at least one final P136 claim in the cached end-to-end run.

\begin{figure}[t]
\centering
\footnotesize
\fbox{\parbox{0.92\linewidth}{\textbf{MARC source}\\Record 001: 990001273190205171\\245: \texthebrew{סדור מנהג אפגניסטן}\\500: \texthebrew{כולל פיוט מאת פנחס חכירי}\\655: \texttt{Piyyutim}}}

\vspace{0.35em}

$\Downarrow$

\vspace{0.35em}

\fbox{\parbox{0.92\linewidth}{\textbf{$X$ fed to the model}\\One text string built from title + local 500 context: \texthebrew{סדור מנהג אפגניסטן ... כולל פיוט מאת פנחס חכירי}}}

\vspace{0.35em}

$\Downarrow$

\vspace{0.35em}

\fbox{\parbox{0.92\linewidth}{\textbf{$Y$ created for training}\\Multi-label vector: \texttt{Piyyutim = 1}; all other genre bits = 0}}
\caption{Model~4 (Genre classifier): distant-supervision example. MARC~655 supplies the label, while MARC~245 + a three-sentence MARC~500 context window form the input text $X$.}
\label{fig:genre-real-example}
\end{figure}

Table~\ref{tab:genre-classifier} summarises the 8 retained genre classes and their training set sizes.

\begin{table}[t]
\caption{Genre classifier training set: 8 retained classes (min.\ 100 examples) after quality filtering.}
\label{tab:genre-classifier}
\centering
\small
\begin{tabular}{lr}
\toprule
\textbf{Genre} & \textbf{N} \\
\midrule
Piyyutim & 5,852 \\
Poetry & 5,464 \\
Illustrated works (Manuscript) & 5,062 \\
Personal correspondence & 3,419 \\
Censored manuscripts & 2,364 \\
Autograph manuscripts & 2,122 \\
Records (Documents) & 1,052 \\
Bibliographies & 775 \\
\midrule
NOTA (abstention class) & 1,629 \\
\midrule
\textbf{Total} & \textbf{25,421} \\
\bottomrule
\end{tabular}
\end{table}

On the pinned subset, this strategy expands genre coverage from 28/68 records with native MARC~655 genre data to 39/68 records with at least one final P136 claim.

% Strategy G (MARC 500 sentence classifier) was removed 2026-05-23.
% Empirical motivation: the Gemini-judged evaluation at confidence
% threshold 0.90 on the pinned 68-record subset showed only 6\%
% strict precision (76\% of high-confidence predictions failed),
% making the classifier head unfit for the auto-approve gate.
% The colophon field is now populated verbatim from MARC; ownership
% extraction remains Provenance NER's responsibility on MARC~561.

\begin{table*}[t]
\caption{MARC-derived training examples for all four models. Companion SVG source files with the same examples are included in the repository.}
\label{tab:real-model-examples}
\centering
\footnotesize
\begin{tabularx}{\textwidth}{>{\raggedright\arraybackslash}p{2.2cm} >{\raggedright\arraybackslash}p{2.2cm} >{\raggedright\arraybackslash}X >{\raggedright\arraybackslash}X >{\raggedright\arraybackslash}X}
\toprule
\textbf{Model} & \textbf{Record 001} & \textbf{MARC fragment} & \textbf{$X$ fed to the model} & \textbf{$Y$ created for training} \\
\midrule
Person NER & 990001253330205171 & \texthebrew{... תפסיר י' דברות מאת אלעזר בן אליעזר} (MARC~500) & Token sequence from the note segment containing \texthebrew{אלעזר בן אליעזר}. & BIO sequence marking \texthebrew{אלעזר בן אליעזר} as \texttt{B-PERSON / I-PERSON / I-PERSON}. \\
\addlinespace
Provenance NER & 990001253170205171 & \texthebrew{ציון בעלים: "מקנת כספי הצעיר יעקב בכ"ר שלמה הכהן"} (MARC~561) & Token sequence from the MARC~561 segment. & BIO sequence marking the ownership phrase as OWNER. \\
\addlinespace
Contents NER & 990001253170205171 & \texthebrew{דף 1א-176ב: ויטל, חיים בן יוסף: אוצרות חיים} (MARC~505) & Token sequence from the reconstructed 505 entry. & BIO sequence for one FOLIO span, one WORK\_AUTHOR span, and one WORK span. \\
\addlinespace
Genre classifier & 990001273190205171 & \texthebrew{סדור מנהג אפגניסטן}; \texthebrew{כולל פיוט מאת פנחס חכירי} (MARC~245 + 500) with MARC~655 = \texttt{Piyyutim} & One text string: title plus the local three-sentence MARC~500 context window. & Multi-label genre vector with \texttt{Piyyutim=1}. \\
\bottomrule
\end{tabularx}
\end{table*}

\begin{table*}[t]
\caption{Model-level contribution on the pinned 68-record verification subset (\texttt{data/tsvs/test\_subset.tsv}).}
\label{tab:model-gap-subset}
\centering
\footnotesize
\begin{tabularx}{\textwidth}{>{\raggedright\arraybackslash}p{2.1cm} >{\raggedright\arraybackslash}X >{\raggedright\arraybackslash}p{2.5cm} >{\raggedright\arraybackslash}p{3.8cm} >{\raggedright\arraybackslash}X}
\toprule
\textbf{Model} & \textbf{Why needed} & \textbf{Potentially fillable gap on pinned subset} & \textbf{Observed Stage~2 output on pinned subset} & \textbf{Observed downstream effect in cached end-to-end subset run} \\
\midrule
Person NER & Finds people mentioned only in free-text notes, especially embedded authors, scribes, translators, and commentators absent from structured headings. & 68/68 records contain note text. & 59 spans across 36 records: 45~AUTHOR, 8~TRANSCRIBER, 4~TRANSLATOR, 1~COMMENTATOR, 1~OWNER. & Feeds the person claim paths; in the cached subset run authorship is asserted as P50 on the 114 created work items (never directly on manuscripts, per the Property:P50 constraint), with P11603 on 1/68. \\
\addlinespace
Provenance NER & Turns prose ownership history into owner/date/collection spans usable for P127 and related claims. & 36/68 records contain direct MARC~561 provenance text. & 13 entities across 15 records: 12~OWNER, 4~DATE, 1~COLLECTION. & The cached subset run yields P127 on 14/68 manuscripts. \\
\addlinespace
Contents NER & Splits one contents string into folios, work titles, and work authors so the pipeline can create work items and exemplar-of links. & 2/68 records contain MARC~505 contents notes. & 93 entities on those two records: 56~FOLIO, 35~WORK, 2~WORK\_AUTHOR. & Supplies structured work evidence behind 126~P1574 claims in the cached subset run. \\
\addlinespace
Genre classifier & Fills P136 when MARC~655 is missing, using title plus local note context instead of direct MARC headings. & 40/68 records have no raw genre headings. & 30 predictions across 28 records; 11 of those records had no raw MARC~655 genre. & Expands the subset from 28 records with native genre data to 39 records with at least one final P136 claim. \\
\bottomrule
\end{tabularx}
\end{table*}


\subsection{Strategy H: Hebrew-to-Latin Transliteration Waterfall}
\label{sec:translit-waterfall}

When a work item or person item has only a Hebrew-script label, Wikidata needs an English label for global discoverability. Naïve algorithmic transliteration of unvocalized Hebrew produces consonant skeletons (e.g.~\code{"Tknot rvno grshm mor hgolh"} for \code{תקנות רבנו גרשם מאור הגולה}) that the Wikidata community rejected as unfit for public labels. We addressed this with a five-tier waterfall that prioritises authoritative sources over algorithmic output and falls through to a stable identifier when no readable transliteration is available.

\paragraph{Five tiers.}
\begin{enumerate}
\item \textbf{Curated override dictionary} (deterministic, $\sim$33 entries): canonical English forms for famous figures and works \textemdash{} \textit{Maimonides}, \textit{Rashi}, \textit{Nahmanides}, \textit{Joseph Karo}, \textit{Ibn Ezra}, \textit{Mishnah}, \textit{Talmud}, etc.

\item \textbf{NLI ALA-LC romanization read} from the source MARC record. NLI catalogers populate field 246 (varying form of title) and 880 (alternate-script linking) with romanized forms following the Library of Congress ALA-LC scheme. When present, this is always preferred over algorithmic output \textemdash{} it is the librarian-quality answer.

\item \textbf{Wikidata SPARQL reverse-lookup}: query the WDQS public endpoint for an item whose \texttt{rdfs:label@he} matches the input exactly, and return its \texttt{rdfs:label@en}. This yields community-consensus forms (e.g.~\code{משה בן מימון} $\to$ \textit{Maimonides}, not algorithmic \textit{Moshe ben Maimon}). Results are cached on disk with 30-day TTL for hits and 24-hour TTL for misses. The lookup respects an \texttt{MHM\_NO\_NETWORK} environment variable for offline-first behaviour.

\item \textbf{TaatikNet ML transliteration}~\cite{taatiknet}: \texttt{malper/taatiknet} is a ByT5-small seq2seq model fine-tuned on $\sim$15k Hebrew$\leftrightarrow$Latin Wiktionary pairs by Morris Alper. We split the input on whitespace and transliterate each word independently (the model was trained on single-word Wiktionary entries and otherwise collapses multi-word phrases). Combining stress marks (acute accents emitted on stressed vowels) are stripped via NFD normalisation. The model is bundled identically in the macOS and Windows installers (\textasciitilde 1.1\,GB), with cross-platform parity enforced by an integration test class.

\item \textbf{Deterministic consonantal ALA-LC fallback}: 27-entry letter table that drops vowel marks. Used only for the person-item P2093 alias path; the work-label path explicitly disables this tier because its output is too ugly for public Wikidata labels.
\end{enumerate}

\paragraph{Compound label format.}
When TaatikNet succeeds, the final English label combines the transliteration with the NLI control number: \code{"Takanut rivno gereshem meor hagola (NLI 990000827290205171)"}. The transliteration is readable; the parenthetical identifier is unambiguous. When TaatikNet returns no usable output (model absent, all-Latin input, etc.), the label collapses to \code{"NLI 990000827290205171"} alone \textemdash{} a stable, unique reference back to the source record. The label is omitted entirely only when neither a transliteration nor a control number is available.

\paragraph{Why distant supervision is the right framing.}
This is the same distant-supervision pattern that powers Strategies A and F: the MARC catalog already carries paired Hebrew$\leftrightarrow$Latin forms in fields 246 and 880, and an existing pretrained model (TaatikNet) was trained on the closest analogous resource (Wiktionary). The waterfall composes these existing sources, falling back to ML only when no librarian-curated answer exists, and to a deterministic identifier when even the ML output is missing. No new annotation effort was required.

\paragraph{Coverage outcome.}
The waterfall eliminates the synthetic \code{"work from Hebrew manuscript X"} placeholder previously flagged by the Wikidata community. On the pinned 68-record subset, all work-level English labels are now either a curated override (Tier~1), an NLI romanization read (Tier~2), a Wikidata-consensus form (Tier~3), or a TaatikNet transliteration paired with the NLI identifier (Tier~4). The fallback to identifier-only labels is itself an upgrade: a Wikidata consumer can resolve the manuscript control number to its NLI record in one click, whereas a synthetic placeholder offered no path forward.

\paragraph{Work items in QuickStatements export.}
The 114 automatically created work items are emitted as their own \code{CREATE} blocks in the QuickStatements dry-run file. The exporter partitions items by \code{entity\_type} and writes them in strict dependency order---persons first (work items reference them via P50), works second (manuscripts reference them via P1574), and manuscripts last---so every back-reference resolved through \code{LAST} points at an entity already declared earlier in the same batch. Without this ordering the work items would still be in \texttt{wikidata\_items.json} but never reach the QuickStatements pipeline, leaving the 126 P1574 claims pointing at unresolved local identifiers.

%% ====================================================================
\section{Data Loss Analysis}
\label{sec:data-loss}
%% ====================================================================

A critical contribution of this work is the systematic quantification of data loss across the MARC-to-Wikidata conversion pipeline. Of 59~fields extracted from MARC records, 37~now map to Wikidata properties (up from 30 after targeted property research). We reduced the data loss rate from 26.4\% to approximately 17\% and identify five remaining categories of loss (Table~\ref{tab:data-loss}).

\begin{table*}[t]
\caption{Taxonomy of data loss in MARC-to-Wikidata conversion on the pinned 68-record subset.}
\label{tab:data-loss}
\centering
\small
\begin{tabular}{p{3.5cm}p{4cm}rp{5.5cm}}
\toprule
\textbf{Category} & \textbf{Fields} & \textbf{Records} & \textbf{Explanation} \\
\midrule
\textbf{1. QID required,}\newline only text available
  & contents (work titles), condition\_notes, related\_works
  & 68, 2, 2
  & Wikidata properties P1574, P5816, P527 require item QIDs; plain work-title strings cannot be asserted directly. \\
\addlinespace
\textbf{2. No Wikidata}\newline property exists
  & has\_multiple\_hands
  & 3
  & No standard Wikidata property for multi-scribe detection. \\
\addlinespace
\textbf{3. Unstructured}\newline free text
  & notes
  & 68
  & General catalog notes with no structured target property. \\
\addlinespace
\textbf{4. Needs further}\newline entity resolution
  & 16 subject terms
  & 16
  & LCSH terms without QID mappings (person names misclassified as subjects). \\
\addlinespace
\textbf{5. Pipeline-internal}\newline metadata
  & attribution\_sources, hierarchy\_type, data\_from\_colophon
  & 68, 68, 25
  & Provenance tracking not intended for Wikidata. \\
\bottomrule
\end{tabular}
\end{table*}

\textbf{The dominant loss remains Category~1}: P1574 (exemplar~of) requires item QIDs, so extracted work-title strings cannot be transferred to Wikidata as literal text. In the pinned subset run, the contents pipeline supplies structured work evidence behind 126~P1574 claims, and the item builder creates 114~work items typed as Q47461344 with title, language, and author when available from NER. This resolves the largest automatically recoverable portion of the loss, while unreconciled or ambiguous work titles remain a structural limitation rather than a parsing failure.

We reduced data loss from 26.4\% to approximately 17\% by identifying Wikidata properties for previously unmapped fields: P2635 (number of codicological units), P1684 (colophon text as inscription with Q372474 qualifier), P7153 (significant places via KIMA resolution), P3132 (explicit/last line), and P478 (volume info); on the pinned 68-record subset these contribute 2, 4, 6, and 2 claims respectively (P478 has no qualifying record in the subset). The counts are small individually, but each mapping removes an entire category of previously untransferable data. Scribal interventions are mapped to P1684 with role qualifiers (Q860740 gloss, Q3299332 correction, Q1136474 marginalia).

Categories~2 and~3 represent \emph{irreducible structural loss}: general catalog notes have no structured Wikidata target, and multi-scribe detection lacks a standard property. Category~4 contains 16~remaining LCSH subject terms (mostly person names misclassified as subjects) awaiting QID mapping.


%% ====================================================================
\section{Experimental Evaluation}
\label{sec:evaluation}
%% ====================================================================

\subsection{Dataset}

The project uses three distinct data scales. The broad source collection is the NLI Hebrew manuscript MARC catalog (123,621~records), which provides the population from which training and evaluation material is drawn. Model training then uses task-specific distant-supervision corpora: for example, the genre classifier scans the full Ktiv export and keeps 25,421 quality-filtered records whose MARC~655 genre headings can be aligned with Hebrew keyword evidence in MARC~500 notes. End-to-end pipeline evaluation, however, is deliberately pinned to \texttt{data/tsvs/test\_subset.tsv}: a deterministic 68-record verification subset rather than a random sample of the full catalog.

The subset was built by \texttt{scripts/build\_test\_subset.py} from two source corpora. The primary source is \texttt{data/tsvs/17th\_century\_samples.tsv} (897 records, SHA256 \texttt{0b0d455...}), which supplies most records and reflects the paper's early-modern manuscript focus. A secondary corpus, \texttt{data/tsvs/filtered\_manuscripts\_after\_906a.tsv} (SHA256 \texttt{ae69b2...}), contributes eight supplemental records for signals not present in the primary corpus. The output subset is pinned by SHA256 \texttt{3f2cd2bb...}; the verification harness refuses to run against a different corpus, making paper claims reproducible across machines.

\begin{table*}[t]
\caption{Deterministic construction of the pinned evaluation subset (\texttt{test\_subset.tsv}).}
\label{tab:test-subset-selection}
\centering
\footnotesize
\begin{tabularx}{\textwidth}{lrlX}
\toprule
\textbf{Stage} & \textbf{N} & \textbf{Selection logic} & \textbf{Purpose} \\
\midrule
1 & 8 & Hard-signal force include & Guarantees rare phenomena such as pre-1582 dates, Hebrew-century dates, Greek/Ladino evidence, MARC~880, and long/multi-owner MARC~561. \\
2 & 9 & Greedy set cover & Adds records covering the largest number of still-uncovered MARC-field signals. \\
3 & 43 & Stratified typical fill & Adds typical records while preserving variation in provenance notes, contributors, and date decades. \\
4 & 8 & Secondary-corpus supplement & Adds edge cases absent from the primary corpus, including reproduction, publication-history, series, and long contents notes. \\
\midrule
\textbf{Total} & \textbf{68} &  & Covers 74 deterministic MARC-field signals used by the test-subset manifest. \\
\bottomrule
\end{tabularx}
\end{table*}

This design separates exploratory data analysis from evaluation. The broad 123,621-record catalog and the 897-record seventeenth-century corpus reveal where information tends to reside in MARC; the pinned 68-record subset is the source of truth for the paper's end-to-end coverage, ablation, and model-contribution numbers.

\subsection{NER Performance}

Table~\ref{tab:ner-results} reports current model performance. All models use DictaBERT~\cite{shmidman2023dictabert} (110M parameters) as the base encoder. The person row reports the v3 role-aware benchmark on the representative held-out sample; provenance, contents, and genre retain the 5-fold validation protocol used for their training runs.

\begin{table}[t]
\caption{Model performance. Person NER is strict exact span+role on the representative held-out benchmark sample; provenance, contents, and genre report their validation protocols.}
\label{tab:ner-results}
\centering
\footnotesize
\setlength{\tabcolsep}{3.5pt}
\begin{tabular}{lrrrl}
\toprule
\textbf{Model} & \textbf{P} & \textbf{R} & \textbf{F1} & \textbf{Training} \\
\midrule
Person NER       & 79.3 & 81.5 & 80.4  & v3 held-out  \\
Provenance NER   & 94.0 & 96.0 & 95.9  & 40min (M4 Pro)  \\
Contents NER     & 100  & 100  & 99.99 & 15min (M4 Pro)  \\
Genre classifier & —    & —    & \textbf{0.918} $\pm$ 0.003  & $\sim$2h (M4 Pro) \\
\bottomrule
\end{tabular}
\end{table}

The person model's strict score is deliberately conservative: a prediction is counted correct only when both the name span and the role match the gold tuple exactly. A second view of the same run separates the two sub-problems: name-only F1 is 86.8\%, and role accuracy among matched names is 92.6\%. Table~\ref{tab:person-v3-gemini} compares the trained v3 model with Gemini extraction under the same exact matching protocol.

\begin{table}[t]
\caption{Person NER v3 benchmark on the representative 100-record sample. Strict F1 counts exact \texttt{(name, role)} tuple matches against gold labels.}
\label{tab:person-v3-gemini}
\centering
\small
\begin{tabular}{lrrrr}
\toprule
\textbf{Method} & \textbf{TP} & \textbf{FP} & \textbf{FN} & \textbf{Strict F1} \\
\midrule
Trained v3 model & 88 & 23 & 20 & 80.4\% \\
Gemini 0-shot    & 67 & 90 & 41 & 50.6\% \\
Gemini 3-shot    & 68 & 86 & 40 & 51.9\% \\
\bottomrule
\end{tabular}
\end{table}

The eval-agent is useful for qualitative error triage but is not the metric used in Table~\ref{tab:person-v3-gemini}. Its verdicts are candidate-level plausibility judgements: they ask whether a predicted candidate \emph{looks right} in the MARC context, but they do not count all missing gold entities as false negatives. On the same run the eval-agent marked 70/111 trained-model candidates as fully correct (63.1\%), while accepting 109/157 Gemini 0-shot candidates (69.4\%) and 108/154 Gemini 3-shot candidates (70.1\%). Because Gemini over-generates candidates and because the judge prompt can accept plausible false positives, strict gold F1 remains the authoritative comparison for model accuracy.

The provenance model's high F1 validates the distant supervision approach: 12,100~training samples generated automatically from MARC field alignment, with 28.4\% augmented with multi-entity examples, yielding a 1.95 percentage point improvement over single-entity training. The contents model's near-perfect F1 reflects the highly regular structure of MARC~505 fields. The genre classifier achieves micro-F1 of 0.918 $\pm$ 0.003 (best fold: 0.921) on 8 classes via a combination of quality-filtered distant supervision (whole-token Hebrew keyword matching, min.\ 100 examples per class), focal loss with per-class positive weighting, and partial fine-tuning of the top 2 DictaBERT layers warm-started from the provenance NER checkpoint. Table~\ref{tab:genre-perclass} reports per-class precision, recall, and F1 averaged across the 5 folds; the weakest class is \emph{Records (Documents)} (F1 0.796), reflecting its smaller training set (1,052 examples) and greater lexical diversity.

These numbers, however, are best read as a \emph{ceiling}, not as a measure of end-to-end contribution. The 5-fold protocol scores each model in isolation against gold spans that were themselves derived (by distant supervision) from MARC structured fields---so it confirms that the models recover the same entities that already exist in those fields. The end-to-end question---``how much new information do the NER models add, beyond what authority linking and the rule layer already produce?''---is answered in \S\ref{sec:evaluation} (End-to-End Coverage) and \S\ref{sec:data-loss} (Data Loss Analysis): on the structured majority of records the marginal contribution is near zero, and the value of the NER models is concentrated in the long tail of records whose entities only ever appear in free-text notes. We retain them in the pipeline as a safety net for exactly that tail.

\begin{table}[t]
\caption{Genre classifier per-class results (mean across 5 folds, threshold = 0.65).}
\label{tab:genre-perclass}
\centering
\small
\begin{tabular}{lrrr}
\toprule
\textbf{Genre} & \textbf{P} & \textbf{R} & \textbf{F1} \\
\midrule
Autograph manuscripts      & 0.994 & 0.942 & 0.967 \\
Personal correspondence    & 0.972 & 0.950 & 0.961 \\
Piyyutim                   & 0.972 & 0.900 & 0.934 \\
Censored manuscripts       & 0.952 & 0.918 & 0.935 \\
Poetry                     & 0.971 & 0.875 & 0.920 \\
Illustrated works          & 0.982 & 0.847 & 0.910 \\
Bibliographies             & 0.908 & 0.854 & 0.880 \\
Records (Documents)        & 0.776 & 0.821 & 0.796 \\
\midrule
NOTA (abstention)          & 0.753 & 0.879 & 0.811 \\
\midrule
\textbf{Micro-F1}          &       &       & \textbf{0.918} \\
\bottomrule
\end{tabular}
\end{table}

\subsection{End-to-End Coverage}

Table~\ref{tab:coverage} reports property-level coverage on the pinned 68-record verification subset.

\begin{table}[t]
\caption{Property-level coverage comparison. P1574 (exemplar of) links each manuscript to its work items, which carry P50 authorship.}
\label{tab:coverage}
\centering
\footnotesize
\setlength{\tabcolsep}{3.5pt}
\begin{tabular}{lrrl}
\toprule
\textbf{Property} & \textbf{Rule} & \textbf{MHM} & \textbf{Strategy} \\
\midrule
P31 (type)       & 100\% & 100\%        & Rule \\
P1476 (title)    & 100\% & 100\%        & Rule \\
P17 (country)    & 0\%   & 100\%        & Hardcoded \\
P131 (located)   & 0\%   & 100\%        & Hardcoded \\
P571 (date)      & 22\%  & \textbf{97\%}& Hebrew parsing \\
P1574 (exemplar of) & 0\% & \textbf{21\%}& Work items \\
P136 (genre)     & 0\%   & \textbf{57\%}& QID mapping + ML \\
P2635 (parts)    & 0\%   & \textbf{3\%}& CU count \\
P921 (subject)   & 0\%   & \textbf{76\%}& Canonical+LCSH \\
P7153 (place)    & 0\%   & \textbf{7\%}& KIMA \\
P1071 (location) & 0\%   & \textbf{7\%}& KIMA \\
P1684 (colophon) & 0\%   & \textbf{6\%}& Text extraction \\
P127 (owner)     & 0\%   & \textbf{21\%}& Provenance NER \\
P11603 (scribe)  & 0\%   & \textbf{1\%}& Person NER \\
\midrule
\textbf{Avg stmts/MS} & \textbf{$\sim$8} & \textbf{30.7} & Combined \\
\bottomrule
\end{tabular}
\end{table}

The multi-strategy approach yields a 3.8$\times$ improvement over rule-only conversion (30.7 vs.\ 8.0 stmts/MS). The dominant share of this gain comes from \emph{authority linking}, which converts every matched person, organisation and place into an internationally-identified entity (Mazal, VIAF, Wikidata, KIMA) and---via VIAF cluster harvesting---multiplies the identifier count $1.7\times$ per matched person, while Wikidata date back-fill (P569/P570) enables the date guardian to vet SPARQL-only candidates. Rule-based parsing handles the Hebrew-specific date formats that authority sources cannot reconstruct (P571: 22\%$\to$97\%). The NER models contribute primarily on the long tail of records whose entities live only in free-text notes: on the structured majority their marginal contribution \emph{beyond} authority + rules is near zero on the pinned evaluation subset, which is why we present them as a safety net rather than as the primary engine.

\subsection{Upload Results}

A build of 248~Wikidata-ready items (Table~\ref{tab:upload}) supports the pipeline's operational feasibility. The count is lower than earlier drafts reported (349) because the 2026-06 authority-hardening pass---homonym abstention and a fail-closed no-identifier validator gate---deliberately declines to build person items that cannot be grounded in at least one external identifier, trading person-item quantity for precision.

\begin{table}[t]
\caption{Wikidata-ready item build on the pinned 68-record evaluation subset.}
\label{tab:upload}
\centering
\small
\begin{tabular}{lr}
\toprule
\textbf{Metric} & \textbf{Value} \\
\midrule
Total items built      & 248 \\
\quad Works            & 114 \\
\quad Persons          & 66 \\
\quad Manuscripts      & 68 \\
Avg stmts/MS           & 30.7 \\
Upload failures        & 0 \\
Duplicate claims       & 0   \\
\bottomrule
\end{tabular}
\end{table}

The zero-failure rate is achieved through property-level type validation (Strategy~E) and claim-level diffing using WikibaseIntegrator's \texttt{MERGE\_REFS\_OR\_APPEND} mode, which adds new claims without duplicating or deleting existing ones. Re-running the upload on previously processed records correctly identifies all 248~items as unchanged, demonstrating idempotency.

\subsection{Ablation Study}

Table~\ref{tab:ablation} shows the contribution of each pipeline component.

\begin{table}[t]
\caption{Component contribution summary on the pinned subset.}
\label{tab:ablation}
\centering
\small
\begin{tabularx}{\linewidth}{p{3.4cm}X}
\toprule
\textbf{Component} & \textbf{Measured contribution} \\
\midrule
Rule-only baseline & Produces approximately 8.0 statements per manuscript, retaining only information already explicit in structured MARC fields. \\
Full pipeline & Produces 30.7 statements per manuscript after combining rules, NER, authority resolution, classifier outputs, and work-item construction. \\
Hebrew date parsing & Raises P571 coverage from the direct structured baseline to 97\% by interpreting MARC~008 plus Hebrew-century and partial-date forms. \\
Person NER & Emits 59 person-role spans across 36/68 records and feeds P50 on the created work items (via the P1574 exemplar chain) plus P11603 on 1/68. \\
Genre classifier & Expands genre coverage from 28/68 records with native MARC~655 data to 39/68 records with final P136 claims. \\
Contents NER + work items & Supplies structured evidence behind 126~P1574 exemplar-of claims and 114 automatically created work items. \\
VIAF cluster harvesting & Adds cross-library identifiers to already matched persons, increasing person-entity richness without changing manuscript count. \\
\bottomrule
\end{tabularx}
\end{table}

The ablation pattern shows that no single component explains the improvement. Rules recover structured dates; NER unlocks people, owners, works, and folio spans from free text; classifiers route otherwise hidden MARC~500 and genre evidence; and authority linking turns local strings into reusable Linked Open Data identifiers.

\subsection{Agentic Evaluation Harness}
\label{sec:eval-harness}

The marginal-contribution and coverage figures above are derived from a deterministic, claim-by-claim comparison against the source MARC record. To evaluate the \emph{semantic} correctness of individual predictions---whether a given NER span, classifier label, or authority match is the right answer for the entity it describes, not merely whether it reproduces a structured field---we developed a separate, large-language-model-as-judge evaluation harness. The harness is a sibling tool rather than a pipeline component: it reads the pipeline's persisted JSON outputs (\texttt{marc\_extracted.json}, \texttt{ner\_results.json}, \texttt{authority\_enriched.json}) from disk and never imports pipeline code or writes back into it, so the system under evaluation and the evaluator remain at arm's length.

The harness was initially a flat workflow: one judge call per candidate prediction. We have since restructured it as a \emph{gated tool-augmented loop} in the ReAct style~\cite{yao2023react}. A cheap first-tier model judges every candidate in a single pass; only candidates whose verdict is uncertain---an explicit \emph{abstain} or a \emph{partially correct}---enter an agentic loop. Inside the loop the model is given a small set of tools that let it gather the evidence a human reviewer would consult before committing to a verdict: \texttt{fetch\_marc\_field} retrieves the raw value of a specific MARC tag, \texttt{expand\_note} returns the full text of a truncated MARC~500 note, \texttt{list\_record\_entities} enumerates the other entities extracted from the same record, and \texttt{lookup\_authority} resolves a candidate against VIAF or Wikidata. If the loop remains uncertain after consulting these tools, the candidate is escalated once to a stronger model for a final verdict. A \texttt{-{}-linear} flag preserves the original single-shot path for reproducibility, and a \texttt{-{}-agentic-all} flag routes every candidate through the loop regardless of first-tier confidence; both the first-tier and the escalation model are user-configurable in the application's settings. This gated design concentrates expensive multi-step reasoning on the small fraction of genuinely ambiguous predictions while keeping the dominant majority of confident verdicts cheap and reproducible.

The harness evaluates two pipeline stages. The original target is Stage~2 (NER and the genre classifier), where each evaluator judges whether an extracted span is a true entity of its declared type. The harness now also evaluates Stage~3 authority resolution: a dedicated \emph{authority} evaluator reads \texttt{authority\_enriched.json} and judges, for each match, whether the resolved authority record (Mazal, VIAF, Wikidata, or KIMA) is in fact the correct real-world entity for the catalogued name---catching the class of error in which a plausible-but-wrong authority record is linked. Authority-stage verification is triggered from a \emph{Verify with AI agent} action on the Stage~3 review surface, mirroring the same action on the Stage~2 surface, so the curator can request an independent second opinion on either extraction or linking before approving.

Beyond binary verdicts, the eval-agent now supports \emph{AI-suggested extraction fixes} for the three NER evaluators (person, provenance, and contents). When the judge determines that a predicted span is partially correct---for example, a person name that omits a patronymic clearly present in the MARC context---it may include a \texttt{suggested\_fix} field in its verdict JSON (schema version~2). A suggested fix carries the corrected text, an optional reasoning string, the source MARC field consulted, and a confidence level; only fixes at the \texttt{high} confidence level are surfaced to the curator. The eval-agent enforces several safety guards before emitting a fix: the corrected text must differ from the original, be non-empty and at most 512~characters, and the evaluator must have access to MARC context or a grounding signal (derived from the entity's \texttt{exists\_in} provenance status: \emph{grounded}~$\to$~\texttt{True}; \emph{wrong\_field} or \emph{novel}~$\to$~\texttt{False}) to ground the correction. Genre-classifier verdicts are explicitly excluded from the fix pathway, as genre labels are not free-text entity strings.

On the curator-facing workbench, a high-confidence suggested fix surfaces as an \textbf{Auto-fix} button attached to the relevant row in the entity table and in the entity detail drawer, alongside a side-by-side preview showing the current and proposed text. Clicking the button sends a \texttt{PATCH} request that updates \texttt{override\_text} without setting \texttt{approved}: the text is corrected but the approval decision remains the curator's. This design ensures the eval-agent can propose a correction while the workbench invariant---\emph{only a human can mark a claim as approved}---is preserved. The suggested fix and its provenance are persisted alongside the verdict in the \texttt{ai\_verdict} JSONB column of the \texttt{extraction\_approvals} table and round-trip through the inference cache, so curators reviewing the same run in a later session see the same proposed correction.

\subsection{Shared Human-in-the-Loop Approval Model}
\label{sec:shared-approval}

A methodological consequence of having two assistive surfaces over the same data---the manual NER and authority editors on one side, the AI-assisted verdict review on the other---is the risk of \emph{divergent} approval state: a curator who approves an entity while editing it manually should not have to re-approve it when reviewing the agent's verdict, and vice versa. We resolve this with a single shared approval store. A \texttt{approvals.json} sidecar is the sole source of truth for which extracted entities and authority matches the curator has accepted; every entry is addressed by a canonical key composed of the manuscript control number, the entity group, the role or type, and the normalised entity text, so the same real-world decision resolves to the same key regardless of which surface recorded it. The NER editor, the authority editor, and the agentic-verdict review UI all read and write this one file, kept consistent across surfaces by a live file watcher, so an approval made in any surface is immediately reflected in the others. This supports a \emph{review once, trust everywhere} workflow: the curator commits to a decision a single time, and that decision is honoured uniformly by the downstream Wikidata-export gate---which continues, unchanged, to project only approved entities---rather than being re-litigated per surface. The shared store thus makes the human's judgement, not any individual UI, the unit of trust in the pipeline.


%% ====================================================================
\section{Discussion}
\label{sec:discussion}
%% ====================================================================

\subsection{Why 100\% Coverage is Structurally Impossible}

Our data loss taxonomy (Section~\ref{sec:data-loss}) demonstrates that the gap between MARC and Wikidata is not merely a technical challenge but a structural one. Three fundamental barriers exist:

\begin{enumerate}[leftmargin=*]
\item \textbf{Entity gap}: Wikidata properties like P1574 (exemplar~of) and P527 (has~parts) require item QIDs. In the pinned subset, the pipeline creates 114 work items and asserts 126 P1574 claims where there is enough structured evidence; remaining ambiguous work strings still require reconciliation or community review before they can be asserted safely.
\item \textbf{Property gap}: Manuscript-specific concepts (colophons, scribal interventions, codicological unit structure) have no Wikidata properties and may never receive them---they represent domain knowledge more suited to specialized ontologies.
\item \textbf{Model gap}: Wikidata's entity-centric model cannot represent relational structures inherent in codicological description, such as ``folios~1--50 contain Work~A in Hand~1, folios~51--100 contain Work~B in Hand~2.''
\end{enumerate}

These barriers motivate a three-layer publication strategy. The HMO RDF graph remains the canonical scholarly representation; a project Wikibase can expose the same HMO-level entities in a browsable and editable knowledge-base interface; and Wikidata receives the richest safe public projection for broad discovery and cross-collection linking. The projection is therefore explicit rather than silent: each HMO class is classified as directly projected to Wikidata, summarized in Wikidata, or retained only in HMO / project Wikibase. The project-Wikibase layer is not a looser duplicate of Wikidata; it is a controlled deployment of the ontology-native graph, intended to retain HMO entities and relationships that public Wikidata should not flatten or over-model.

\subsection{The Distant Supervision Advantage}

Our distant supervision methodology has a key practical advantage: it requires zero manual annotation. The MARC dual structure---structured fields as labels, note fields as training text---is an inherent property of cataloging workflows worldwide. Any institution with MARC records containing both controlled headings and descriptive notes can replicate our approach for their collection, regardless of language or subject domain.

The methodology's effectiveness varies with cataloging quality: NLI's records, created by professional catalogers over decades, provide high-quality distant labels. Less rigorously cataloged collections may require additional noise handling.

\subsection{VIAF Cluster Harvesting as LOD Amplification}

The cluster harvesting technique demonstrates that LOD enrichment is multiplicative: a single successful authority match (VIAF) cascades into 2--4 additional cross-references (GND, LCCN, ISNI, BnF), each connecting the entity to a different national library's knowledge graph. For the 61~persons with VIAF matches in the pinned subset build (92\% of all built persons), cluster harvesting added 102~additional authority identifiers---an amplification factor of 1.7$\times$.

This technique is generalizable to any domain where VIAF covers the entities of interest. The rate-limiting constraint (2~requests/second for search + 1 additional for cluster fetch) means processing 123K~records requires approximately 17~hours for the authority resolution stage, which is acceptable for batch processing.

\subsection{Domain Expert Validation}

The pipeline's Wikidata output was reviewed by two domain experts in Hebrew manuscript studies from the University of Haifa (M.~Lavee and E.~Baumgarten, October 2025). Their feedback, grounded in hands-on cataloging of seven test manuscripts (Vatican~44, Huntington~115, Oppenheimer~129, Parma~3122, two London manuscripts, and Jerusalem~8210), identified four gaps in the Wikidata representation:

\begin{enumerate}[leftmargin=*]
\item \textbf{Certainty/confidence mechanism}: No formal way to express that a scribe attribution is ``probable'' vs.\ ``certain''. The experts requested a vocabulary equivalent to \emph{certain / probable / possible} on individual claims.
\item \textbf{Partial work identification}: Manuscripts frequently contain compositions that partially overlap or derive from known works without being identical to them. The experts requested relations analogous to \texttt{is\_variant\_of} or \texttt{possibly\_realises} (IFLA LRM terminology) to express this grey area.
\item \textbf{Multi-volume and anthology structure}: The BU--CU--PU hierarchy is insufficiently expressive for manuscripts where one codicological unit contains nested sub-units, or where multiple volumes collectively form a single work.
\item \textbf{Textual overlap between scribal hands}: No formal relation for cases like Huntington~115, where two scribes copied overlapping sections of the same text.
\end{enumerate}

We addressed items (1) and (2) directly: unconfirmed person attributions (those without an authority-matched Wikidata QID) now receive a \texttt{P1480~=~Q18122778} (\emph{presumably}) qualifier; unreconciled work links (P1574) receive the same qualifier, implementing the \texttt{possibly\_realises} semantics within Wikidata's property model. Items (3) and (4) remain structural limitations of Wikidata (see below), so they are preserved in the HMO RDF graph and in the project-Wikibase export rather than flattened into public Wikidata claims.

\subsection{Community Validation and Safety Guards}

A pilot upload of 670~items to Wikidata in April~2026 generated substantive community feedback through user talk pages. The Wikidata community identified: (a)~duplicate person items caused by the reconciler checking only VIAF and NLI identifiers, missing existing items findable via LCCN/GND/ISNI; (b)~institutional holders (MARC~710 entries) incorrectly attached as P50~(author) instead of P195~(collection) or P127~(owned~by); (c)~P8189 authority-record identifiers attached to non-human entities; and (d)~generic catalog placeholders (Hebrew \textit{kovetz}, `collection') emitted verbatim as Wikidata labels.

Each issue was fixed and protected by a regression test. The full reconciler now checks all five identifier types before creating any new person entity; the institutional name detector routes MARC~710 entries to the appropriate collection/ownership property; P8189 is gated on both authority-record prefix and \texttt{P31=Q5} (human); and placeholder titles are detected and replaced with shelfmark-based labels. Fourteen pre-upload validation checks in \texttt{item\_validator.py} block any item carrying a known failure pattern (empty label, anonymous person, kovetz placeholder, trailing punctuation, MARC-inverted label, institutional-as-person, P3959 on human, P8189 bad prefix, VIAF on non-person, multiple same-language P1559, no-identifier person, P1559 Latin-as-Hebrew, Hebrew label is Latin, ambiguous single name) from being approved in the GUI; every error-level finding disables the ``Approved'' checkbox for that row.

Beyond data-quality checks, a secondary class of incident arose when an automated cleanup script modified items the pipeline had never created. To make this class of error structurally impossible, the uploader now enforces a \textbf{four-stage modification guard} (\texttt{CLAUDE.md} rule~38) gated by \textbf{three independent verification channels}. The four gates are: (1)~an entry-point check in \code{upload\_item()}; (2)~a re-assertion inside \code{\_build\_wbi\_item()} so new call paths cannot bypass gate~1; (3)~a per-statement identity-conflict check refusing any value that would create a multi-value conflict on an existing identifier; and (4)~a last-ditch check placed immediately before the single \code{wbi\_item.write()} call site in the code-base. Raising \code{UnauthorisedModificationError} at any gate converts the attempt into a clean \emph{skipped} result rather than a partial edit. The three verification channels, all cross-checked inside \code{\_is\_our\_item()}: (i)~the MediaWiki API's \code{prop=revisions\&rvdir=newer} endpoint for the authoritative first-revision author; (ii)~the MediaWiki API's \code{list=usercontribs\&uctype=new} endpoint as an independent confirmation that the authenticated user has a \emph{page-creation} edit on the QID; and (iii)~a Wikidata SPARQL \code{ASK} query confirming the QID still exists and has not been deleted, redirected, or blanked since reconciliation. A negative answer from any channel refuses the edit; a network failure on one channel does not unlock the gate as long as the primary revisions answer still agrees.

The 545~unit tests in the safety/regression suite serve as a barrier ensuring these classes of error cannot recur. Three of those tests are \emph{structural} rather than behavioural --- they fail if a future refactor introduces a second \code{wbi\_item.write()} call, separates the pre-write guard from the write, or re-introduces a marker-based creator fallback --- so the architecture of Rule~38 cannot be unintentionally weakened.

\subsection{Interactive Validation Workbench}

To make the full safety and validation architecture accessible to domain curators without CLI access, the pipeline exposes a \emph{Wikidata Studio} workbench that unifies the review/check/approve/upload workflow into one surface. Given a corpus of ~5\,785~pipeline-built items per 100~manuscripts, three design choices drive scalability and reviewability:

\textbf{Validate-with-Wikidata pass.} A single \code{\_ValidationWorker} thread visits every item and performs all live Wikidata checks in one pass: reconciler duplicate detection, per-identifier collision checking (for each of P214, P8189, P244, P227, P213, querying Wikidata for the current owner of the identifier and flagging any disagreement with our item's \code{existing\_qid}), SPARQL \code{ASK} existence checking on every matched QID (catching items deleted, redirected, or blanked between reconciliation and approval), three-channel creator verification per item, and a fresh re-run of \code{validate\_item()} on the possibly-curator-edited item. A full payload---including the actual SPARQL ASK string fired, per-PID verdict matrix, and the three-channel creator result---is emitted per row so the curator has full accountability for each status decision.

\textbf{Pagination.} A chained \code{\_PaginationProxy} truncates the viewport's \code{rowCount()} to the configured page size (25 / 50 / 100), so the \code{QTableView} materialises at most 100 rows regardless of backing corpus size. Row-level action widgets (edit, claims, detail-view) attach only to currently-visible rows. Empirically, this reduces scroll latency for 5\,785-row corpora from $\sim$1.2\,s to unmeasurable and keeps widget count bounded at $< 300$ (pagination bar + visible-page actions).

\textbf{Persistent validation state.} After validation completes, the workbench exports a \code{<input\_stem>\_wikidata\_verified.json} file that round-trips through the load path. Each entry carries the full \code{WikidataItem} (labels, descriptions, aliases, statements with qualifiers and references) plus the validation payload (status, reason, matched QID, identifier checks, SPARQL ask/answer, creator check, validator issues, timestamp, latency) and the approved flag. The curator can close the application, reopen the file later, and resume review from the exact state they left, including any per-item approval decisions already made. This is essential for large corpora where review may span multiple sessions.

\textbf{RDF-first source and projection audit.} The workbench now treats \texttt{output.ttl} as the preferred Stage~6 input. When a reviewed authority file exists beside it, the projection reuses that sidecar only for curator-approved authority decisions; the HMO RDF graph remains the source being projected. Every RDF-based run emits three audit artefacts: \texttt{wikidata\_source\_report.json} records whether the source was raw authority data, reviewed authority data, HMO RDF, or a saved Wikidata-review state; \texttt{wikidata\_projection\_coverage.json} reports how each HMO class is represented or intentionally not represented in Wikidata; and \texttt{wikibase\_entities.json} / \texttt{wikibase\_quickstatements.txt} provide an offline project-Wikibase draft for the full HMO graph. In the current implementation, the project-Wikibase draft is complete at the entity/statement level relative to the generated HMO graph: every typed HMO node is represented as a draft entity, and all non-type/non-label RDF assertions are carried as statement drafts. Multi-typed RDF resources currently expose one preferred class in the JSON draft; adding an explicit list of all RDF classes is a planned schema-fidelity refinement before live project-Wikibase upload.

Every interactive validation surface (status-detail popup, claim editor, credential entry, full-screen browser) renders against a shared Liquid-Glass \code{GraphBackdrop} (CLAUDE.md rule~37) so the workbench visually reads as a single application rather than a stack of pop-up dialogs. The design aligns with the Apple HIG~2026 Liquid-Glass aesthetic.

\textbf{Authority-match biodata comparison.} Reviewing authority matches (Mazal, VIAF, KIMA) against the source MARC record is a second bottleneck: each match carries a label, a source, and an ID, but the curator's approval decision hinges on biographical data --- dates, places, occupations, name variants in multiple scripts --- that is not surfaced in the row summary. We expose a per-row \emph{biodata-comparison dialog} (opened by a \textbf{🧬} icon next to the existing edit/source icons) with five tabs: \emph{Dates}, \emph{Places}, \emph{Names / variants}, \emph{Occupations}, and \emph{Raw}. All four data sources (MARC record, Mazal authority row, VIAF cluster JSON, KIMA place entry) flow through pure extractors that emit a common \code{BioData} dataclass, so the dialog renders side-by-side diff pills coloured by verdict --- \emph{matched}, \emph{differs}, \emph{only-in-MARC}, \emph{only-in-authority} --- after Unicode NFKC + case-fold normalisation. The MARC extractor is entity-scoped: given the specific row being reviewed, it retrieves only the fields belonging to the matched entity (author/contributor/place), using a token-overlap matcher that strips Hebrew name particles ``בן''/``בת'' and Latin particles ``de''/``di''/``ibn'' plus single-letter initials to avoid false positives across people with common patronymics. The \emph{Raw} tab shows both sides as pretty-printed JSON with cross-key value matches ($\ge 0.90$ similarity) highlighted in blue --- matches are detected on extracted scalar values regardless of which key they live under (so MARC's \code{occupations:["author"]} and an authority's \code{notes:["role: author"]} both light up), while purely structural JSON lines (empty objects, closing brackets) are excluded to prevent trivial matches from dominating the signal. VIAF cluster fetches run on \code{QThreadPool} with a 0.5s rate-limit gate; results cache process-wide by \code{(source, authority\_id)}; the MARC side renders synchronously on open so the dialog is usable before the authority fetch completes. Approve + Next buttons in the footer allow the curator to walk through all matches without dismissing the dialog, turning what was a 20-modal click-through into keyboard-paced review.

\textbf{Auth-match table with field-origin Source column.} The Stage~3 review surface emits one row per data-bearing entry in the record: every MARC main, added, subject, or series entry (100/110/111/700/710/711 and, after v1.10, 600/610/651/752/800/810/811), every NER entity from the three models, and every KIMA place lookup --- matched or not. Unmatched NER entities (the OWNER, DATE, WORK, FOLIO, and WORK\_AUTHOR rows that survive auto-approval but lack authority IDs) are kept visible so the curator reviews them in the same surface as the MARC matches, rather than only the ones that happened to resolve. The \emph{Source} column shows the field or pipeline of origin --- \texttt{MARC~100} for personal main entries, \texttt{MARC~710} for corporate added entries, \texttt{MARC~600} for subject persons, \texttt{MARC~800} for series added persons, \texttt{NER Author} for person-NER spans, \texttt{NER Owner} for provenance-NER OWNER spans, \texttt{KIMA Place} for Hebrew place lookups, and similarly for the other field/pipeline labels --- rather than the authority service that resolved the match (which the curator reads from the \texttt{matched\_id} and \texttt{wikidata\_qid} columns instead). This decoupling lets the curator filter by data origin (``show me every \texttt{NER Owner} row regardless of authority outcome'') and authority outcome (``show me every row with no Wikidata QID'') as independent axes, supporting the three-axis Source~$\times$~Type~$\times$~Role filter announced in §1.3.

\textbf{Confidence-column hover tooltips.} The tri-level Stage~3 verdict (\emph{high} / \emph{medium} / \emph{low}) is a compact signal but hides the reasoning that produced it. Approving or rejecting a borderline match was previously a black-box decision: the curator could see the colour band but not the underlying signal. We now render a structured HTML tooltip on hover over the \emph{Conf.} column that enumerates every signal that fed the verdict --- which authority sources matched (Mazal / VIAF / Wikidata / KIMA), whether a Latin preferred name was found for cross-script verification, the biographical years that drove the date-conflict guard, any guard flag that fired (\texttt{date\_conflict}, \texttt{short\_name\_homonym}, \texttt{cluster\_collapse}, \texttt{over\_merge\_detected}, \texttt{wikidata\_confirms}, \texttt{cross\_source\_conflict}) with a one-line human-readable description, the harvested VIAF cluster identifiers (GND, LCCN, ISNI, BnF), and, when a guard hard-rejected the match, the explicit \texttt{rejection\_reason} string (e.g.\ ``date-conflict: person born 1700, MS dated 1500''). A companion tooltip on the NER editor's confidence columns distinguishes the keyword-classifier bucket (0.60 / 0.85, the signal that Stage~3 guards key on) from the BIO transformer softmax averaged across the entity's tokens (the real model probability), so the curator can write auto-approve rules against the appropriate signal. Both tooltips are pure HTML rendered via Qt's \code{QToolTip} so the theme colours stick on macOS, and they read from underscore-prefixed signal keys carried by the flattened row dict; the flattening layer was extended in v1.10 to copy \texttt{guard\_flags}, \texttt{rejection\_reason}, \texttt{sources}, \texttt{source\_count}, \texttt{birth\_year}, \texttt{death\_year}, and \texttt{preferred\_name\_lat} from \texttt{authority\_enriched.json} so the tooltip can render without re-running Stage~3.

\subsection{Limitations}

\begin{enumerate}[leftmargin=*]
\item \textbf{Test set bias}: The 68-record verification subset is intentionally coverage-oriented rather than random. It includes rare and difficult MARC phenomena, so it is excellent for regression testing and claim verification, but average coverage on the full catalog may differ.
\item \textbf{Person NER evaluation scope}: The v3 role-aware model fixes the earlier keyword-role workaround, but its 80.4\% strict span+role F1 is measured on a representative 100-record held-out benchmark sample rather than the full 123K-record catalog. The sample is intentionally useful for regression and comparison, but larger random and curator-reviewed evaluations are still needed before treating the number as a catalog-wide estimate.
\item \textbf{VIAF rate limiting}: At 2~requests/second, full authority resolution of 123K~records requires significant processing time. Caching and batch scheduling mitigate but do not eliminate this constraint.
\item \textbf{Language specificity}: The Hebrew date parser, NER models, and keyword heuristics are Hebrew-specific. Adapting to other languages requires retraining NER models and creating new parsing rules, though the distant supervision methodology transfers directly.
\item \textbf{Wikidata codicological expressivity}: Manuscript-specific relational structures (nested scribal layers, textual overlap, variant-work relationships at sub-folio granularity) exceed what Wikidata's property model can represent. The pipeline therefore treats Wikidata as a public summary projection, while the full HMO RDF graph --- and, optionally, a project-Wikibase deployment of that graph --- preserve the richer scholarly structure.
\end{enumerate}


%% ====================================================================
\section{Conclusion}
\label{sec:conclusion}
%% ====================================================================

We presented the MHM Pipeline as three contributions that share a single processing pipeline. \emph{First}, authority enrichment across thirteen MARC tags (100/110/111, 600/610, 651, 700/710/711, 752, 800/810/811) against four complementary sources (Mazal/NLI, VIAF, Wikidata, KIMA), amplified by VIAF cluster harvesting and Wikidata date back-fill, is the primary engine of new information per manuscript---it turns every matched entity from a string into a linked, internationally-identified item and multiplies the international identifier count $1.7\times$ per matched person. \emph{Second}, the same pipeline publishes its output in three coordinated Linked-Open-Data layers: the full HMO research graph as RDF, a project Wikibase at \texttt{mhm-hmo.wikibase.cloud} for browsing and curator editing, and a public Wikidata projection that reaches 87\% of HMO classes and averages 30.7~statements per manuscript on the pinned 68-record verification subset. \emph{Third}, Wikidata Studio puts a researcher-facing workbench around the whole flow: a Source~$\times$~Type~$\times$~Role review filter, a side-by-side biographic-comparison dialog, a confidence tooltip that surfaces every guardian and cluster identifier behind a match, 14 pre-upload guards, and resumable JSON review state---so every statement that reaches Wikidata is one a human approved with the evidence in view.

The distant-supervision methodology that exploits MARC's dual structure---structured fields as automatic labels for the same information in free-text notes---is a separable methodological contribution that yields three NER models. The current role-aware person model reaches 80.4\% strict span+role F1 on a representative held-out benchmark sample (86.8\% name-only F1; 92.6\% role accuracy among matched names), while provenance and contents NER reach 95.9\% and 99.99\% F1 respectively. End-to-end, however, their marginal contribution \emph{beyond} authority enrichment and rules is near zero on the structured majority of records on the pinned evaluation subset, and we retain them in the pipeline as a safety net for the long tail of records whose entities live only in free-text notes. This is a deliberate honesty about where the value actually comes from: authority linking and the workbench are doing the work; the NER models cover the cases authority cannot reach.

Automatic work item creation supplies 114~work items and 126~P1574 exemplar-of claims, addressing the largest automatically recoverable portion of the structural loss identified in our taxonomy. Our data loss taxonomy identifies five categories of information that cannot transfer cleanly to Wikidata, demonstrating that the dominant remaining limitation is structural rather than technical---an argument for the three-layer publication model we adopt rather than for a single Wikidata-only target. A verified build of 248~Wikidata-ready items (68 manuscripts, 66 identifier-grounded persons, 114 works) achieved zero failures and zero duplicate claims, demonstrating the feasibility of the guarded upload workflow.

For the cultural heritage community, the take-away is that a high-coverage MARC-to-LOD pipeline is best framed as an \emph{authority enrichment + multi-layer publication + curator workbench} system, with statistical extraction (NER, classifiers) used to cover the tail rather than as the headline. All code, models, training data generation scripts and the Wikidata Studio application are available under GPL-3.0 at \url{https://github.com/alexandergolbergwix/pipeline}.

\subsection{Future Work}

Five directions are planned: (1)~expanding work-item reconciliation beyond the 114 automatically created subset items so more P1574 claims can be confirmed rather than qualified; (2)~evaluating the pipeline on the full 123,621-record NLI catalog; (3)~adapting the distant supervision methodology to Arabic and Latin manuscript collections to test cross-lingual transferability; (4)~replacing the \emph{presumably} qualifier on NER-identified works with confirmed Wikidata reconciliation once the work-item corpus is established; and (5)~deploying the project-Wikibase layer for the full HMO graph, including explicit preservation of multi-class RDF typing in the project-Wikibase draft schema, so codicological units, textual witnesses, text traditions, and epistemological structures can be reviewed and queried without forcing them into Wikidata's public property model.


%% ====================================================================
%% REFERENCES
%% ====================================================================

\begin{thebibliography}{99}

\bibitem{loc-marc}
Library of Congress.
\newblock MARC~21 Format for Bibliographic Data.
\newblock \url{https://www.loc.gov/marc/bibliographic/}, 2024.

\bibitem{vrandecic2014wikidata}
Vrande\v{c}i\'{c}, D. and Kr\"{o}tzsch, M.
\newblock Wikidata: A Free Collaborative Knowledge Base.
\newblock \emph{Communications of the ACM}, 57(10):78--85, 2014.

\bibitem{van2020wikidata}
van Veen, T.
\newblock Wikidata: From ``an'' Identifier to ``the'' Identifier.
\newblock \emph{Information Technology and Libraries}, 38(2):72--81, 2020.

\bibitem{mcallister2017bibframe}
McCallum, S.
\newblock BIBFRAME Development.
\newblock \emph{JLIS.it}, 8(3):71--85, 2017.

\bibitem{cornell-ld4l}
Cornell University Library.
\newblock Linked Data for Libraries (LD4L).
\newblock \url{https://ld4l.org/}, 2016.

\bibitem{ehrmann2023named}
Ehrmann, M., Hamdi, A., Pontes, E.L., Romanello, M., and Doucet, A.
\newblock Named Entity Recognition and Classification in Historical Documents: A Survey.
\newblock \emph{ACM Computing Surveys}, 56(2):1--47, 2023.

\bibitem{ehrmann2022extended}
Ehrmann, M., Romanello, M., Fl\"{u}ckiger, A., and Clematide, S.
\newblock Extended Overview of HIPE-2022: Named Entity Recognition and Linking in Multilingual Historical Documents.
\newblock In \emph{CLEF 2022}, pp.~1038--1066, 2022.

\bibitem{shmidman2023dictabert}
Shmidman, A., Shmidman, S., Koppel, M., and Goldberg, Y.
\newblock DictaBERT: A State-of-the-Art BERT Suite for Modern Hebrew.
\newblock \emph{arXiv preprint arXiv:2304.09919}, 2023.

\bibitem{goldberg2025person}
Goldberg, A., Prebor, G., and Elmalech, A.
\newblock Person NER for Hebrew Manuscript Catalogs via Distant Supervision.
\newblock \emph{Submitted}, 2025.

\bibitem{mintz2009distant}
Mintz, M., Bills, S., Snow, R., and Jurafsky, D.
\newblock Distant Supervision for Relation Extraction without Labeled Data.
\newblock In \emph{ACL-IJCNLP 2009}, pp.~1003--1011, 2009.

\bibitem{ratner2017snorkel}
Ratner, A., Bach, S.H., Ehrenberg, H., Fries, J., Wu, S., and R\'{e}, C.
\newblock Snorkel: Rapid Training Data Creation with Weak Supervision.
\newblock In \emph{VLDB}, 11(3):269--282, 2017.

\bibitem{lin2017focal}
Lin, T.-Y., Goyal, P., Girshick, R., He, K., and Doll\'{a}r, P.
\newblock Focal Loss for Dense Object Detection.
\newblock In \emph{IEEE International Conference on Computer Vision (ICCV)}, pp.~2980--2988, 2017.

\bibitem{yao2023react}
Yao, S., Zhao, J., Yu, D., Du, N., Shafran, I., Narasimhan, K., and Cao, Y.
\newblock ReAct: Synergizing Reasoning and Acting in Language Models.
\newblock In \emph{International Conference on Learning Representations (ICLR)}, 2023.

\bibitem{feitosa2023wikidata}
Feitosa, D., Dermaku, G., and Ioannidis, L.
\newblock Wikidata for Digital Libraries: A Systematic Literature Review.
\newblock \emph{Journal of Documentation}, 79(7):1--25, 2023.

\bibitem{marcedit}
Reese, T.
\newblock MarcEdit: Your Complete Free MARC Editing Utility --- Features.
\newblock \url{https://marcedit.reeset.net/features}, 2024.

\bibitem{openrefine-recon}
OpenRefine.
\newblock Reconciling --- OpenRefine User Manual.
\newblock \url{https://openrefine.org/docs/manual/reconciling}, 2024.

\bibitem{openrefine-wikibase}
OpenRefine.
\newblock Overview of Wikibase Support --- OpenRefine User Manual.
\newblock \url{https://openrefine.org/docs/manual/wikibase/overview}, 2024.

\bibitem{aalberg2013}
Aalberg, T. and \v{Z}umer, M.
\newblock The Value of MARC Data, or, Challenges of Frbrisation.
\newblock \emph{Journal of Documentation}, 69(6):851--872, 2013.

\bibitem{taatiknet}
Alper, M.
\newblock TaatikNet: Sequence-to-Sequence Learning for Hebrew Transliteration.
\newblock Model: \url{https://huggingface.co/malper/taatiknet}. Source: \url{https://github.com/morrisalp/taatiknet}, 2023.

\end{thebibliography}

\end{document}
